%%%%%%%%%%%%%%%%%%%%%%%%%%%%%%%%%%%%%%%%%%%%%%%%%%%%%%%%%%%%%%%%%%%%%%
% Lớp 10 | Đại số tổ hợp | Nhị thức Newton
% Dạng 0D8?3-4: Tìm hệ  số, số hạng trong khai triển
% Số câu: 84  (sinh tự động bởi tach_cau.py; từ nay thầy sửa trực tiếp file này)
%%%%%%%%%%%%%%%%%%%%%%%%%%%%%%%%%%%%%%%%%%%%%%%%%%%%%%%%%%%%%%%%%%%%%%

%%%%%%%%%%%%%%%% TRẮC NGHIỆM 4 PHƯƠNG ÁN %%%%%%%%%%%%%%%%

% nguồn: MaHoaNH_TracNghiem_Lop10
\begin{ex}%[0D8N3-4]%[ma:81408668]
	Số hạng chính giữa trong khai triển $(3 x+2 y)^4$ là 
	\choice 
		{$\mathrm{C}_{4}^2 x^2  y^2$}
		{$(3 x)^2(2 y)^2$} 
		{$6 \mathrm{C}_{4}^2  x^2 y^2$}
		{\True $36  \mathrm{C}_{4}^2 x^2 y^2$}
	\loigiai{ 
		Ta có $(3 x+2 y)^4=(3 x)^4+4(3  x)^3(2 y)+6(3 x)^2(2 y)^2+4(3 x)(2 y)^3+(2 y)^4$.\\
		Suy ra hệ số chính giữa trong khai triển trên là $6(3 x)^2(2  y)^2=36 \mathrm{C}_{4}^2 x^2 y^2$.
	 }
\end{ex}

% nguồn: MaHoaNH_TracNghiem_Lop10
\begin{ex}%[0D8N3-4]%[ma:b7b29c53]
	Trong khai triển $(2 a-b)^5$, hệ số của số hạng thứ $3$  bằng
	 \choice
		{$-80$} 
		{\True $80$} 
		 {$-10$}
		{$10$} 
	 \loigiai{
		Ta  có
		\begin{eqnarray*}
			(2  a-b)^5 &=&(2 a)^5-5(2 a)^4 b+10(2 a)^3 b^2-10(2 a)^2 b^3+5(2 a) b^4-b^5\\
			&=&32 a^5-80 a^4 b+80 a^3 b^2-40 a^2 b^3+10 a b^4-b^5. 
		\end{eqnarray*}
	} 
\end{ex}

% nguồn: MaHoaNH_TracNghiem_Lop10
\begin{ex}%[0D8N3-4]%[ma:d1cf9fcc]
	Tìm hệ số  của $a^3 b^2$ trong khai triển nhị thức $(a+2 b)^5$.
	\choice 
		 {$160$}
		{$80$} 
		{$20$} 
		 {\True $40$}
	 \loigiai{
		Ta  có
		\begin{eqnarray*}
			(a+2 b)^5&=&a^5+5 a^4(2 b)+10 a^3(2 b)^2+10 a^2(2 b)^3+5 a(2 b)^4+(2  b)^5\\
			&=&a^5+10 a^4  b+40 a^3 b^2+80 a^2 b^3+80 a b^4+32 b^5
		\end{eqnarray*}
		Vậy hệ số của $a^3 b^2$ trong khai triển là $40$. 
	} 
\end{ex}

% nguồn: MaHoaNH_TracNghiem_Lop10
\begin{ex}%[0D8H3-4]%[ma:641ad6d0]
	Trong khai triển $(0{,}3 + 0{,}7)^4$, số  hạng thứ hai là?
	\choice 
		{\True $0{,}0756$} 
		 {$0{,}2646$}
		 {$0{,}3752$}
		{$0{,}2463$} 
	 \loigiai{
		Trong khai triển $(0{,}3 + 0{,}7)^4$, số hạng thứ  hai là $\mathrm{C}_{4}^{1}\cdot 0{,}3^{3}\cdot 0{,}7^1 = 0{,}0756$.
	 }
\end{ex}

% nguồn: MaHoaNH_TracNghiem_Lop10
\begin{ex}%[0D8H3-4]%[ma:f9e3258d]
	Tìm số hạng chứa $x^2$ trong khai triển  $P(x)=4x^2+x(x-2)^4$.
	\choice 
		 {$28x^2$}
		 {\True $-28x^2$}
		{$-24x^2$} 
		 {$24x^2$}
	\loigiai{ 
		Ta có $P(x)=4 x^2+x(x-2)^4=4 x^2+x \displaystyle\sum\limits_{k=0}^4 \mathrm{C}_{4}^k x^{4-k}(-2)^k=4 x^2+\displaystyle\sum\limits_{k=0}^4 \mathrm{C}_{4}^k(-2)^k x^{5-k}$.\\ 
		Số  hạng chứa $x^2$ (ứng với $k=3$) trong khai triển $P(x)$ là $\left[4+\mathrm{C}_{4}^3(-2)^3\right] x^2=-28x^2$.
	} 
\end{ex}

% nguồn: MaHoaNH_TracNghiem_Lop10
\begin{ex}%[0D8H3-4]%[ma:91d6244a]
	Trong khai triển $(2x+1)^5$ hệ số của số hạng  chứa $x^5$ là
	\choice 
		{\True $32$} 
		 {$10$}
		 {$100$}
		 {$1\,000$}
	 \loigiai{
		Ta có $(2x+1)^5=\displaystyle\sum\limits_{i=0}^5 \mathrm{C}_{5}^{i} (2x)^{5-i} 1^{i}=\displaystyle\sum\limits_{i=0}^5 \mathrm{C}_{5}^{i} 2^{5-i}  x^{5-i}$.\\
		Do đó hệ số của số hạng  chứa $x^5$ trong khai triển là $\mathrm{C}_{5}^0 2^5=32$.
	} 
\end{ex}

% nguồn: MaHoaNH_TracNghiem_Lop10
\begin{ex}%[0D8H3-4]%[ma:5714c9d7]
	 Trong khai triển $(1+3x)^4$, số hạng thứ $2$ theo số mũ tăng dần của $x$ là
	 \choice
		 {$108x$}
		 {$54x^2$}
		 {$1$}
		{\True $12x$} 
	 \loigiai{
		Ta có $(1+3x)^4=\displaystyle\sum\limits_{k=0}^4 \mathrm{C}_{4}^k(3x)^k=\displaystyle\sum\limits_{k=0}^4 \mathrm{C}_{4}^k 3^k  x^k$.\\
		Do đó số hạng thứ $2$ theo số mũ  tăng dần của $x$ ứng với $k=1$, tức là $\mathrm{C}_{4}^1 3^1 x=12x$.
	 }
\end{ex}

% nguồn: MaHoaNH_TracNghiem_Lop10
\begin{ex}%[0D8H3-4]%[ma:c2ceabe9]
	Dùng hai số hạng đầu tiên trong khai triển $(x+\Delta x)^4$  để tính gần đúng số $(1{,}01)^4$. Tìm số đó.
	\choice 
		{\True  $1{,}04$}
		 {$1{,}0406$}
		 {$1{,}040604$}
		 {$1{,}04060401$}
	\loigiai{ 
		 $(1{,}01)^4=(1+0{,}01)^4=\mathrm{C}_4^0+\mathrm{C}_4^1\cdot 0{,}01+\mathrm{C}_4^2\cdot 0{,}01^2+\mathrm{C}_4^3\cdot 0{,}01^3+\mathrm{C}_4^4\cdot 0{,}01^4$.\\
		 Khi đó $(1{,}01)^4\approx \mathrm{C}_4^0+\mathrm{C}_4^1\cdot 0{,}01=1{,}04$.
	} 
\end{ex}

% nguồn: MaHoaNH_TracNghiem_Lop10
\begin{ex}%[0D8H3-4]%[ma:f966c84b]
	Số hạng thứ $4$ trong khai triển Newton của  đa thức $P(x)=(1+x)^{18}$ là
	 \choice
		{$\mathrm{C}_{18}^4$} 
		{$\mathrm{C}_{18}^4x^4$} 
		 {$\mathrm{C}_{18}^3$}
		{\True  $\mathrm{C}_{18}^3x^3$}
	\loigiai{ 
		Số hạng tổng quát trong khai triển của $P(x)$ là $\mathrm{C}_{18}^k1^{18-k}x^k$.\\ Số  hạng thứ tư ứng với $k=3$.\\
		Do đó số hạng thứ tư là  $\mathrm C_{18}^31^{15}x^3=\mathrm{C}_{18}^3x^3$.
	} 
\end{ex}

% nguồn: MaHoaNH_TracNghiem_Lop10
\begin{ex}%[0D8H3-4]%[ma:6b2a8e4c]
	Khi khai triển biểu thức $\left(\sqrt{2}x-1\right)^5$, số hạng chứa  $x^4$ là
	 \choice
		{\True $-20x^4$} 
		{$20\sqrt{2}x^4$} 
		{$-20$} 
		 {$20x^4$}
	\loigiai{ 
		Ta  có $\left(\sqrt{2}x-1\right)^5=\mathrm{C}_5^0\left(\sqrt{2}x\right)^5-\mathrm{C}_5^1\left(\sqrt{2}x\right)^4+\mathrm{C}_5^2\left(\sqrt{2}x\right)^3-\mathrm{C}_5^3\left(\sqrt{2}x\right)^2+\mathrm{C}_5^4\left(\sqrt{2}x\right)-\mathrm{C}_5^5$.\\
		Số hạng chứa $x^4$  trong khai triển là $-\mathrm{C}_5^1(\sqrt{2})^4\cdot x^4=-20x^4$.
	} 
\end{ex}

% nguồn: MaHoaNH_TracNghiem_Lop10
\begin{ex}%[0D8H3-4]%[ma:3b1f1f6f]
	Tìm hệ số của $x^3$ trong khai triển  $(2x+1)^5$.
	 \choice
		 {$-80$}
		{$10$} 
		 {$40$}
		{\True $80$} 
	\loigiai{ 
		Ta có $(2x+1)^5=\mathrm{C}_{5}^0(2x)^5+\mathrm{C}_{5}^1(2x)^4+\mathrm{C}_{5}^2(2x)^3+\mathrm{C}_{5}^3(2x)^2+\mathrm{C}_{5}^4  2x+\mathrm{C}_{5}^5$.\\
		Do đó số hạng chứa $x^3$ là $\mathrm{C}_{5}^3(2x)^3=8  \mathrm{C}_{5}^3 x^3=80x^3$.\\
		Vậy hệ số của $x^3$ trong khai triển  $(2x+1)^5$ là $80$.
	 }
\end{ex}

% nguồn: MaHoaNH_TracNghiem_Lop10
\begin{ex}%[0D8H3-4]%[ma:78871136]
	Trong khai triển $\left(x-\dfrac{2}{x^2}\right)^{12}$ số hạng chứa $x^6$ là 
	 \choice
		{\True  $264 x^6$}
		 {$264$}
		 {$100 x^6$}
		 {$100$}
	 \loigiai{
		Ta có $\left(x-\dfrac{2}{x^2}\right)^{12}=\displaystyle\sum\limits_{i=0}^{12} \mathrm{C}_{12}^{i}  x^{12-i}\left(\dfrac{-2}{x^2}\right)^{i}=\displaystyle\sum\limits_{i=0}^{12} \mathrm{C}_{12}^{i}(-2)^{i} x^{12-3 i}$.\\
		 Từ yêu cầu bài toán ta có $12-3 i=6 \Leftrightarrow i=2$.\\
		Do  đó số hạng chứa $x^6$ trong khai triển là $\mathrm{C}_{12}^2(-2)^2 x^6=264 x^6$.
	} 
\end{ex}

% nguồn: MaHoaNH_TracNghiem_Lop10
\begin{ex}%[0D8H3-4]%[ma:17fab1fc]
	Tìm số hạng không  chứa $x$ trong khai triển $\left(\dfrac{1}{x}+x^3\right)^4$.
	 \choice
		 {$1$}
		{\True  $4$}
		{$6$} 
		{$12$} 
	 \loigiai{
		 Ta có $\left(\dfrac{1}{x}+x^3\right)^4=\displaystyle\sum\limits_{k=0}^4 \mathrm{C}_{4}^k\left(\dfrac{1}{x}\right)^{4-k}\left(x^3\right)^k=\displaystyle\sum\limits_{k=0}^4 \mathrm{C}_{4}^k x^{4 k-4}$.\\
		Số hạng không chứa $x$ trong  khai triển trên ứng với $4 k-4=0 \Leftrightarrow k=1$.\\
		Vậy số hạng không chứa  $x$ trong khai triển $\left(\dfrac{1}{x}+x^3\right)^4$ là $\mathrm{C}_{4}^1=4$.
	} 
\end{ex}

% nguồn: MaHoaNH_TracNghiem_Lop10
\begin{ex}%[0D8H3-4]%[ma:0a0afec7]
	Trong khai triển $(a + b)^5$, hệ số của số hạng thứ  $3$ bằng
	 \choice
		 {$\mathrm{C}_{5}^{3}$}
		 {\True $\mathrm{C}_{5}^{2}$}
		{$\mathrm{C}_{6}^{3}$} 
		 {$\mathrm{C}_{6}^{2}$}
	\loigiai{ 
		Số hạng tổng quát thứ $k+1$ trong khai triển của $(a + b)^5$ là  $\mathrm{C}_{5}^{k}a^{k}b^{5-k}$.\\
		Như vậy, số hạng tổng quát thứ $3$ trong khai triển của $(a + b)^5$ là $\mathrm{C}_{5}^{2}a^{2}b^{5-2}  = \mathrm{C}_{5}^{2}a^{2}b^{3}$.\\
		Do đó, hệ số của số  hạng thứ $3$ bằng $\mathrm{C}_{5}^{2}$.
	 }
\end{ex}

% nguồn: MaHoaNH_TracNghiem_Lop10
\begin{ex}%[0D8H3-4]%[ma:8c28873a]
	Tìm hệ số của số hạng chứa $x^2y^2$ trong khai triển $(x+2y)^{4}$. 
	\choice 
		{\True  $24$}
		 {$16$}
		 {$32$}
		 {$4$}
	\loigiai{ 
		Số hạng tổng quát của khai triển là $\mathrm{C}_{4}^{k}\cdot x^{4-k}\cdot (2y)^{k}=\mathrm{C}_{4}^{k}\cdot2^k\cdot x^{4-k}\cdot y^{k}$ ($0\leq k\leq n$, $k\in\mathbb{N}$).\\ 
		Số hạng này chứa  $x^2y^2$ khi và chỉ khi $\heva{&4-k=2\\&k=2}\Leftrightarrow k=2$.\\
		Vậy hệ số  của số hạng chứa $x^2y^2$ là $\mathrm{C}_{4}^{2}\cdot2^{2}=24$.
	} 
\end{ex}

% nguồn: MaHoaNH_TracNghiem_Lop10
\begin{ex}%[0D8H3-4]%[ma:9797f300]
	Dùng bốn số hạng đầu tiên trong khai triển $(x+\Delta x)^4$ để tính gần đúng số $(4{,}001)^4$. Tìm số  đó.
	 \choice
		 {\True $256{,}2560963$}
		 {$256{,}25$}
		 {$256{,}256$}
		{$256{,}256096$} 
	\loigiai{ 
		$$ 
		(4{,}001)^4=(4+0{,}001)^4=\mathrm{C}_4^0 \cdot 4^4+\mathrm{C}_4^1 \cdot 4^3 \cdot 0{,}001+\mathrm{C}_4^2 \cdot 4^2 \cdot 0{,}001^2+\mathrm{C}_4^3 \cdot 4^3 \cdot 0{,}001^3+\mathrm{C}_4^4 \cdot  4^4 \cdot 0{,}001^4.
		$$ 
		Khi đó $(4{,}001)^4\approx \mathrm{C}_4^0\cdot 4^4+\mathrm{C}_4^1\cdot 4^3\cdot 0{,}001+\mathrm{C}_4^2\cdot 4^2\cdot  0{,}001^2+\mathrm{C}_4^3\cdot 4^3\cdot 0{,}001^3=256{,}2560963$.
	 }
\end{ex}

% nguồn: MaHoaNH_TracNghiem_Lop10
\begin{ex}%[0D8H3-4]%[ma:2bed2cbd]
	Tìm hệ số của $x^4$ trong khai triển  $(2x-3)^5$
	\choice 
		{$-270$} 
		 {$-80$}
		 {$240$}
		{\True  $-240$}
	 \loigiai{
		 Ta có
		$$(2x-3)^5=\mathrm{C}_{5}^0(2x)^5+\mathrm{C}_{5}^1(2x)^4(-3)+\mathrm{C}_{5}^2(2x)^3(-3)^2+\mathrm{C}_{5}^3(2x)^2(-3)^3+\mathrm{C}_{5}^4(2x)(-3)^4+\mathrm{C}_{5}^5(-3)^5.$$ 
		Do đó  số hạng chứa $x^4$ là $\mathrm{C}_{5}^1(2x)^4(-3)=-240x^4$.\\
		Vậy hệ số của  $x^4$ trong khai triển $(2x-3)^5$ là $-240$.
	 }
\end{ex}

% nguồn: MaHoaNH_TracNghiem_Lop10
\begin{ex}%[0D8H3-4]%[ma:0b2bcd17]
	Dùng bốn số hạng đầu tiên trong khai triển $(x+\Delta x)^5$ để tính gần  đúng số $(1{,}03)^5$. Tìm số đó.
	 \choice
		 {$1{,}15$}
		{$1{,}1592$} 
		{\True  $1{,}159274$}
		 {$1{,}15927407$}
	\loigiai{ 
		$$ 
		(1{,}03)^5=(1+0{,}03)^5=\mathrm{C}_5^0+\mathrm{C}_5^1 \cdot 0{,}03+\mathrm{C}_5^2 \cdot 0{,}03^2+\mathrm{C}_5^3 \cdot 0{,}03^3+\mathrm{C}_5^4 \cdot 0{,}03^4+\mathrm{C}_5^5 \cdot 0{,}03^5. 
		 $$
		Khi  đó $(1{,}03)^5\approx \mathrm{C}_5^0+\mathrm{C}_5^1\cdot 0{,}03+\mathrm{C}_5^2\cdot 0{,}03^2+\mathrm{C}_5^3\cdot 0{,}03^3=1{,}159274$.
	} 
\end{ex}

% nguồn: MaHoaNH_TracNghiem_Lop10
\begin{ex}%[0D8H3-4]%[ma:29f47179]
	Dùng ba  số hạng đầu tiên trong khai triển $(x+\Delta x)^4$ để tính gần đúng số $(1{,}02)^4$. Tìm số đó.
	\choice 
		 {$1{,}08$}
		{\True  $1{,}0824$}
		{$1{,}08243$} 
		{$1{,}082432$} 
	\loigiai{ 
		 $$
		(1{,}02)^4=(1+0{,}02)^4=\mathrm{C}_4^0+\mathrm{C}_4^1 \cdot 0{,}02+\mathrm{C}_4^2 \cdot  0{,}02^2+\mathrm{C}_4^3 \cdot 0{,}02^3+\mathrm{C}_4^4 \cdot 0{,}02^4.
		 $$
		Khi đó $(1{,}02)^4\approx  \mathrm{C}_4^0+\mathrm{C}_4^1\cdot 0{,}02+\mathrm{C}_4^2\cdot 0{,}02^2=1{,}0824$.
	} 
\end{ex}

% nguồn: MaHoaNH_TracNghiem_Lop10
\begin{ex}%[0D8H3-4]%[ma:d7514453]
	Dùng ba số hạng đầu tiên trong khai triển $(x+\Delta x)^5$ để tính gần đúng số $(2{,}03)^5$. Tìm số đó. 
	\choice 
		 {\True $34{,}473$}
		{$34{,}47$} 
		{$34{,}47308$} 
		{$34{,}473088$} 
	 \loigiai{
		 $$
		(2{,}03)^5=(2+0{,}03)^5=\mathrm{C}_5^0 \cdot 2^5+\mathrm{C}_5^1 \cdot 2^4  \cdot 0{,}03+\mathrm{C}_5^2 \cdot 2^3 \cdot 0{,}03^2+\mathrm{C}_5^3 \cdot 2^2 \cdot 0{,}03^3+\mathrm{C}_5^4 \cdot 2 \cdot 0{,}03^4+\mathrm{C}_5^5 \cdot 0{,}03^5.
		$$ 
		Khi đó $(2{,}03)^5\approx \mathrm{C}_5^0\cdot 2^5+\mathrm{C}_5^1\cdot 2^4\cdot 0{,}03+\mathrm{C}_5^2\cdot  2^5\cdot 0{,}03^2=34{,}473$.
	 }
\end{ex}

% nguồn: MaHoaNH_TracNghiem_Lop10
\begin{ex}%[0D8H3-4]%[ma:a2a26758]
	Tìm  hệ số của $x^2$ trong khai triển: $f(x)=\left(x^3+\dfrac{2}{x^2}\right)^n$, với $x > 0$, biết tổng ba hệ số đầu của $x$ trong khai triển bằng 33.
	 \choice
		 {$34$}
		{\True  $24$}
		{$6$} 
		{$12$} 
	 \loigiai{
		 Ta có $\mathrm{C}_n^0+2\mathrm{C}_n^1+4\mathrm{C}_n^2=33\Rightarrow n=4$.\\
		Số hạng tổng quát của khai  triển $f(x)=\left(x^3+\dfrac{2}{x^2}\right)^4$ là \[T_{k+1}=\mathrm{C}_4^k(x^3)^{4-k}\left(\dfrac{2}{x^2}\right)^k=2^k \mathrm{C}_4^k x^{12-5k}.\]
		Số hạng chứa $x^2$ trong  khai triển ứng với số mũ của $x$ là $12-5k=2\Leftrightarrow k=2$.\\
		Vậy  hệ số của $x^2$ trong khai triển là $2^2\mathrm{C}_4^2=24$.
	} 
\end{ex}

% nguồn: MaHoaNH_TracNghiem_Lop10
\begin{ex}%[0D8H3-4]%[ma:d80856a6]
	Dùng ba số  hạng đầu tiên trong khai triển $(x+\Delta x)^5$ để tính gần đúng số $(2{,}01)^5$. Tìm số đó.
	 \choice
		 {$32{,}808$}
		 {$32{,}80804$}
		{\True $32{,}8$} 
		{$32{,}8080401$} 
	\loigiai{ 
		 $$
		(2{,}01)^5=(2+0{,}01)^5=\mathrm{C}_5^\circ \cdot 2^5+\mathrm{C}_5^1 \cdot  2^4 \cdot 0{,}01+\mathrm{C}_5^2 \cdot 2^3 \cdot 0{,}01^2+\mathrm{C}_5^3 \cdot 2^2 \cdot 0{,}01^3+\mathrm{C}_5^4 \cdot 2 \cdot 0{,}01^4+\mathrm{C}_5^5 \cdot 0{,}01^5.
		$$ 
		Khi đó  $(2{,}01)^5\approx \mathrm{C}_5^\circ\cdot 2^5+\mathrm{C}_5^1\cdot 2^4\cdot 0{,}01=32{,}8$.
	 }
\end{ex}

% nguồn: MaHoaNH_TracNghiem_Lop10
\begin{ex}%[0D8H3-4]%[ma:b4586465]
	Biết hệ số của $x^2$ trong khai triển của $(1-3x)^n$ là 90. Khi đó ta  có $3n^4$ bằng
	 \choice
		{$7\,203$} 
		{\True $1\,875$} 
		{$1\,296$} 
		{$6\,561$} 
	 \loigiai{
		Số hạng tổng quát  khai triển của $(1-3x)^n$ là $T_{k+1}=\mathrm{C}_n^k(-3x)^k=(-3)^k \mathrm{C}_n^k x^k$.\\
		Suy ra hệ số của $x^2$ trong khai triển của $(1-3x)^n$ ứng với  $k=2$.\\
		Khi đó $(-3)^2\mathrm{C}_n^2=90\Leftrightarrow 9\dfrac{n(n-1)}{2}=90\Leftrightarrow n(n-1)=20\Leftrightarrow\left[\begin{array}{l}n=-4\\ n=5
	\end{array} \Rightarrow 3n^4=1\,875\right.$.
} 
\end{ex}

% nguồn: MaHoaNH_TracNghiem_Lop10
\begin{ex}%[0D8H3-4]%[ma:fe00c24f]
	Trong khai triển biểu thức $(x+7)^{4}$, hệ số của số hạng  chính giữa là
	 \choice
		{$28$} 
		{\True  $294$}
		{$1372$} 
		{$2401$} 
	\loigiai{ 
		Số hạng thứ $k+1$ của khai triển là $T_{k+1}=\mathrm{C}_{4}^{k}\cdot x^{4-k}\cdot 7^{k}$ ($0\leq k\leq  4$, $k\in\mathbb{N}$).\\
		Trong khai triển của  nhị thức $(x+7)^{4}$ có $5$ số hạng.\\
		Suy ra số hạng chính giữa  là số hạng thứ ba, ứng với $k=2$.\\
		Vậy hệ số của số hạng đứng  giữa là $\mathrm{C}_{4}^{2}\cdot7^{2}=294$.
	} 
\end{ex}

% nguồn: MaHoaNH_TracNghiem_Lop10
\begin{ex}%[0D8H3-4]%[ma:01c9e35a]
	 Hệ số của $x^2 y^2$ trong khai triển $(x+2y)^4$ là
	\choice 
		{$32$} 
		{$8$} 
		{\True $24$} 
		{$16$} 
	\loigiai{ 
		Ta có $(x+2y)^4=\displaystyle\sum\limits_{k=0}^4 \mathrm{C}_{4}^k x^{4-k}(2y)^k=\displaystyle\sum\limits_{k=0}^4 \mathrm{C}_{4}^k \cdot 2^k \cdot x^{4-k}  y^k$.\\
		Số hạng chứa $x^2 y^2$ trong khai triển trên  ứng với $\heva{
		&4-k=2  \\
		&k=2}  \Leftrightarrow k=2$.\\
		Vậy  hệ số của $x^2 y^2$ trong khai triển $(x+2y)^4$ là $\mathrm{C}_{4}^2 \cdot 2^2=24$.
	} 
\end{ex}

% nguồn: MaHoaNH_TracNghiem_Lop10
\begin{ex}%[0D8H3-4]%[ma:1c0f5eb9]
	 Tìm hệ số của $x^2$ trong khai triển: $f(x)=\left(x^3+\dfrac{1}{x^2}\right)^n$, với $x > 0$, biết\break $\mathrm{C}_n^0+\mathrm{C}_n^1+\mathrm{C}_n^2=11$.
	 \choice
		{$20$} 
		 {\True $6$}
		 {$7$}
		 {$15$}
	\loigiai{ 
		Ta có $\mathrm{C}_n^0+\mathrm{C}_n^1+\mathrm{C}_n^2=11\Leftrightarrow 1+n+\dfrac{n(n-1)}{2}=11\Leftrightarrow\left[\begin{array}{l}n=4\\ n=-5.
	\end{array}\right.$\\
	Số hạng tổng quát của khai triển $f(x)=\left(x^3+\dfrac{1}{x^2}\right)^4$  là \[T_{k+1}=\mathrm{C}_4^k(x^3)^{4-k}\left(\dfrac{1}{x^2}\right)^k=\mathrm{C}_4^k x^{12-5k}.\]
	Số hạng chứa $x^2$ trong khai triển  ứng với số mũ của $x$ là $12-5k=2\Leftrightarrow k=2$.\\
	Vậy hệ số của $x^2$  trong khai triển là $\mathrm{C}_4^2=6$.
 }
\end{ex}

% nguồn: MaHoaNH_TracNghiem_Lop10
\begin{ex}%[0D8H3-4]%[ma:47035cd0]
	Tìm số hạng không chứa $x$ trong  khai triển nhị thức Newton của $\left(\dfrac{1}{x}+x^3\right)^4$.
	\choice 
		{$1$} 
		{\True  $4$}
		 {$6$}
		{$12$} 
	 \loigiai{
		Khai triển nhị thức ta  có
		\allowdisplaybreaks 
		\begin{eqnarray*}
			\left(\dfrac{1}{x}+x^3\right)^4&=&\mathrm{C}_4^0\left(\dfrac{1}{x}\right)^4+\mathrm{C}_4^1\left(\dfrac{1}{x}\right)^3\cdot  x^3+\mathrm{C}_4^2\left(\dfrac{1}{x}\right)^2(x^3)^2+\mathrm{C}_4^3\left(\dfrac{1}{x}\right)(x^3)^3+\mathrm{C}_4^4(x^3)^4\\
			&=&x^{12}+4x^8+6x^4+4+\dfrac{1}{x^4}. 
		\end{eqnarray*}
		Do đó, số hạng không chứa $x$ trong khai triển là $4$. 
	} 
\end{ex}

% nguồn: MaHoaNH_TracNghiem_Lop10
\begin{ex}%[0D8H3-4]%[ma:3ee145b4]
	Xét khai triển $(2x-3)^4=a_4x^4+a_3x^3+a_2x^2+a_1x+a_0$, giá trị của  $a_2+a_1$ là
	\choice 
		{\True  $0$}
		{$-216$} 
		 {$216$}
		 {$-432$}
	 \loigiai{
		 Ta có
		\allowdisplaybreaks 
		$\begin{aligned}
		 (2x-3)^4&=\mathrm{C}_4^0\cdot(2x)^4+\mathrm{C}_4^1\cdot(2x)^3\cdot(-3)^1+\mathrm{C}_4^2\cdot(2x)^2\cdot(-3)^2+\mathrm{C}_4^3\cdot(2x)^1\cdot(-3)^3+\mathrm{C}_4^4\cdot(-3)^4\\
		&=16x^4-96x^3+216x^2-216x+81. 
	\end{aligned}$\\
	 Vậy $a_2+a_1=216-216=0$.
} 
\end{ex}

% nguồn: MaHoaNH_TracNghiem_Lop10
\begin{ex}%[0D8H3-4]%[ma:bd713f5d]
	Dùng bốn số hạng đầu tiên trong khai triển $(x+\Delta  x)^5$ để tính gần đúng số $(4{,}0002)^5$. Tìm số đó.
	 \choice
		{$1024{,}25$} 
		{\True  $1024{,}256026$}
		 {$1024{,}25602$}
		 {$1024{,}256$}
	 \loigiai{
		$$ 
		\begin{aligned}
			(4{,}0002)^5=(4+0{,}0002)^5=4^5\cdot \mathrm{C}_5^0&+4^4\cdot \mathrm{C}_5^1\cdot 0{,}0002+4^3\cdot \mathrm{C}_5^2\cdot  0{,}0002^2+4^2\mathrm{C}_5^3\cdot 0{,}0002^3
			 \\
			&+4\mathrm{C}_5^4\cdot 0{,}0002^4+\mathrm{C}_5^5\cdot  0{,}0002^5.
		\end{aligned}
		$$ 
		Khi  đó
		$(4{,}0002)^5\approx \mathrm{C}_5^0\cdot 4^5+\mathrm{C}_5^1\cdot 4^4\cdot 0{,}0002+\mathrm{C}_5^2\cdot 4^3\cdot 0{,}0002^2+\mathrm{C}_5^3\cdot  4^2\cdot 0{,}0002^3=1024{,}256026$.
	 }
\end{ex}

% nguồn: MaHoaNH_TracNghiem_Lop10
\begin{ex}%[0D8H3-4]%[ma:18e73e1d]
	Dùng bốn số hạng đầu tiên trong khai triển $(x+\Delta  x)^5$ để tính gần đúng số $(2{,}0003)^5$. Tìm số đó.
	\choice 
		 {$32{,}02$}
		 {$32{,}024$}
		{\True $32{,}0240072$} 
		 {$32{,}024007$}
	 \loigiai{
		 $$
		\begin{aligned}
			(2{,}0003)^5=(2+0{,}0003)^5=2^5 \cdot \mathrm{C}_5^0 &+2^4  \cdot \mathrm{C}_5^1 \cdot 0{,}0003+2^3 \cdot \mathrm{C}_5^2 \cdot 0{,}0003^2+2^2 \mathrm{C}_5^3 \cdot 0{,}0003^3 \\
			&+2 \mathrm{C}_5^4 \cdot  0{,}0003^4+\mathrm{C}_5^5 \cdot 0{,}0003^5.
		\end{aligned}
		$$ 
		Khi đó 
		$$ 
		(2{,}0003)^5 \approx \mathrm{C}_5^0 \cdot 2^5+\mathrm{C}_5^1  \cdot 2^4 \cdot 0{,}0003+\mathrm{C}_5^2 \cdot 2^3 \cdot 0{,}0003^2+\mathrm{C}_5^3 \cdot 2^2 \cdot 0{,}0003^3=32{,}0240072.
		$$ 
	 }
\end{ex}

% nguồn: MaHoaNH_TracNghiem_Lop10
\begin{ex}%[0D8H3-4]%[ma:15f7dff9]
	Giả sử có khai triển  $(1-2x)^n=a_0+a_1 x+a_2 x^2+\cdots+a_n x^n$. Tìm $a_4$ biết $a_0+a_1+a_2=31$.
	\choice 
		{\True  $80$}
		{$-80$} 
		 {$40$}
		 {$-40$}
	 \loigiai{
		Ta  có
		 \allowdisplaybreaks
		\begin{align*}
			(1-2x)^n&=\mathrm{C}_n^01^n(-2x)^0+\mathrm{C}_n^11^{n-1}(-2x)+\mathrm{C}_n^21^{n-2}(-2x)^2+\cdots 
			\\ 
			&=1-2\mathrm{C}_n^1x+4\mathrm{C}_n^2x^2+\cdots 
		\end{align*}
		Suy ra  $a_0=1; a_1=-2\mathrm{C}_n^1; a_2=4\mathrm{C}_n^2$.\\
		Theo bài ra $a_0+a_1+a_2=31$ nên ta  có
		 \allowdisplaybreaks
		\begin{align*}
			 &1-2\mathrm{C}_n^1+4\mathrm{C}_n^2=31
			\\ 
			\Leftrightarrow& 1-2\dfrac{n!}{1!(n-1)!}+4\dfrac{n!}{2!(n-2)!}=31 
			\\ 
			\Leftrightarrow&  1-2n+2n(n-1)=31
			 \\
			 \Leftrightarrow& 2n^2-4n-30=0
			 \\
			 \Leftrightarrow& n^2-2n-15=0
			 \\
			 \Rightarrow& n=5.
		\end{align*}
		Từ  đó ta có $a_4=\mathrm{C}_5^4(-2)^4=80$.
	} 
\end{ex}

% nguồn: MaHoaNH_TracNghiem_Lop10
\begin{ex}%[0D8V3-4]%[ma:4efd93ec]
	Tìm số hạng chứa  $x^2$ trong khai triển nhị thức $P(x)=4x^2+x(x-2)^4$.
	\choice 
		 {$28x^2$}
		{\True $-28x^2$} 
		{$-24x^2$} 
		 {$24x^2$}
	\loigiai{ 
		Ta có  $P(x)=4x^2+x(x-2)^4=4x^2+x\sum\limits_{k=0}^4\mathrm{C}_4^kx^{4-k}(-2)^k
		=4x^2+\sum\limits_{k=0}^4\mathrm{C}_4^k(-2)^kx^{5-k}$.\\ 
		Số hạng chứa $x^2$ (ứng với $k=3$)  trong khai triển $P(x)$ là $\left[4+\mathrm{C}_4^3(-2)^3\right]x^2=-28x^2$.
	} 
\end{ex}

% nguồn: MaHoaNH_TracNghiem_Lop10
\begin{ex}%[0D8V3-4]%[ma:08376d46]
	Trong khai triển biểu thức $F=\left(\sqrt{3}+\sqrt[3]2\right)^5$  số hạng nguyên có giá trị lớn nhất là
	 \choice
		{$8$} 
		 {\True $60$}
		{$58$} 
		 {$20$}
	\loigiai{ 
		Ta có  số hạng tổng quát $T_{k+1}=\mathrm{C}_5^k\left(\sqrt{3})^{5-k}(\sqrt[3]2\right)^k$.\\
		Ta thấy bậc hai của căn thức là $2$ và $3$ là hai số nguyên  tố, do đó để $T_{k+1}$ là một số nguyên
		thì $\left\{\begin{array}{l}k \in \mathbb{N} \\ 0\leq k \leq 5\\ (5-k) \vdots 2\\ k \vdots 3
	\end{array} \Leftrightarrow k=3\Rightarrow T_4=\mathrm{C}_5^3\left(\sqrt{3}\right)^2(\sqrt[3]2)^3\right.$.\\
	Vậy trong khai triển  có giá trị lớn nhất là số hạng nguyên là $T_4=60$.
	 }
\end{ex}

% nguồn: MaHoaNH_TracNghiem_Lop10
\begin{ex}%[0D8V3-4]%[ma:ad63a52b]
	Gọi $n$ là số  nguyên dương thỏa mãn $\mathrm{A}_n^3+2\mathrm{A}_n^2=48$. Tìm hệ số của $x^3$ trong khai triển nhị thức Newtơn của $(1-3x)^n$.
	 \choice
		 {\True $-108$}
		 {$81$}
		 {$54$}
		{$-12$} 
	 \loigiai{
		Điều kiện $n\ge 3$; $n\in \mathbb{N}$. Khi đó,  ta có
		\begin{eqnarray*}
			\mathrm{A}_n^3+2\mathrm{A}_n^2=48 &\Leftrightarrow&  \dfrac{n!}{(n-3)!}+2\cdot \dfrac{n!}{(n-2)!}=48 \\
			&\Leftrightarrow&  n(n-1)(n-2)+2\cdot n(n-1)=48 \\
			&\Leftrightarrow& n^3-n^2-48=0 \Leftrightarrow  n=4 (\text{nhận}).
		\end{eqnarray*}
		Ta  có $(1-3x)^4=\sum\limits_{k=0}^4\mathrm{C}_4^k\left(-3x\right)^k
		 =\sum\limits_{k=0}^4\mathrm{C}_4^k\left(-3\right)^kx^k$.\\
		Hệ số của $x^3$ trong khai triển trên ứng với  $k=3$.\\
		Vậy hệ số của $x^3$ trong khai triển $(1-3x)^4$ là  $\mathrm{C}_4^3\cdot \left(-3\right)^3=-108$.
	} 
\end{ex}

% nguồn: MaHoaNH_TracNghiem_Lop10
\begin{ex}%[0D8V3-4]%[ma:189880b8]
	Gọi $n$ là số nguyên dương thỏa mãn  $\mathrm{A}_n^3+2 \mathrm{A}_n^2=48$. Tìm hệ số của $x^3$ trong khai triển $(1-3 x)^n$.
	 \choice
		{\True $-108$} 
		 {$81$}
		 {$54$}
		 {$-12$}
	 \loigiai{
		Điều kiện  $n \geq 3$, $n \in \mathbb{N}$.\\
		 Ta có
		\begin{eqnarray*}
			\mathrm{A}_n^3+2 \mathrm{A}_n^2=48 &\Leftrightarrow& \dfrac{n!}{(n-3)!}+2 \cdot \dfrac{n!}{(n-2)!}=48\\ 
			&\Leftrightarrow& n(n-1)(n-2)+2 \cdot  n(n-1)=48\\
			&\Leftrightarrow& n^3-n^2-48=0 \Leftrightarrow n=4. 
		\end{eqnarray*}
		 Ta có $(1-3 x)^4=\displaystyle\sum\limits_{k=0}^4 \mathrm{C}_{4}^k(-3 x)^k=\displaystyle\sum\limits_{k=0}^4 \mathrm{C}_{4}^k(-3)^k x^k$.\\
		Hệ số của $x^3$ trong khai triển ứng  với $k=3$.\\
		Vậy hệ  số của $x^3$ trong khai triển $(1-3 x)^4$ là $\mathrm{C}_{4}^3 \cdot(-3)^3=-108$.
	} 
\end{ex}

% nguồn: MaHoaNH_TracNghiem_Lop10
\begin{ex}%[0D8V3-4]%[ma:1e37bbfa]
	Với $n$ là số nguyên dương, gọi $a_{3n-3}$ là hệ số của $x^{3n-3}$ trong khai triển thành đa thức của $f(x)=\left(x^2+1\right)^n(x+2)^n$.  Tìm $n$ để $a_{3n-3}=26n$.
	\choice 
		{$n=11$} 
		{\True $n=5$} 
		{$n=12$} 
		{$n=10$} 
	 \loigiai{
		$$ 
		\begin{aligned}
			& f(x)=\left(x^2+1\right)^n(x+2)^n=\left(\sum\limits_{k=0}^n \mathrm{C}_n^k  x^{2 n-2 k}\right)\left(\sum\limits_{i=0}^n \mathrm{C}_n^i x^{n-i} 2^i\right)=\sum\limits_{k=0}^n\left(\sum\limits_{i=0}^n \mathrm{C}_n^k \mathrm{C}_n^i 2^i x^{3 n-2 k-i}\right), \\
			& (0 \leq  i, k \leq n) \\
			& \text{Yêu cầu} \Leftrightarrow 3 n-(2 k+i)=3 n-3 \Leftrightarrow 2 k+i=3 \Leftrightarrow\left[\begin{array}{l}
			k=i=1 \\ 
			k=0,  i=3.
		\end{array}\right. \\
		\text{Suy ra}~ &a_{3n-3}=2  \mathrm{C}_n^1 \mathrm{C}_n^1+2^3 \mathrm{C}_n^0 \mathrm{C}_n^3=26 n \Leftrightarrow n=5.
	\end{aligned}
	 $$
} 
\end{ex}

% nguồn: MaHoaNH_TracNghiem_Lop10
\begin{ex}%[0D8V3-4]%[ma:60c4961a]
	Cho khai triển $(1+2x)^n=a_0+a_1 x+a_2 x^2+\cdots+a_n  x^n$, biết $n$ thỏa mãn $a_0+8a_1=2a_2+1$. Tìm hệ số lớn nhất của khai triển.
	 \choice
		{$160$} 
		{\True $80$} 
		 {$60$}
		{$105$} 
	\loigiai{ 
		Ta có $(1+2x)^n=a_0+a_1 x+a_2 x^2+\cdots+a_n  x^n=\sum\limits_{k=0}^n \mathrm{C}_n^k(2x)^k=\sum\limits_{k=0}^n \mathrm{C}_n^k 2^k x^k$.\\
		$\Rightarrow a_k=\mathrm{C}_n^k  2^k \Rightarrow a_0=\mathrm{C}_n^0$, $a_1=2\mathrm{C}_n^1$, $a_2=2^2\mathrm{C}_n^2$.\\
		Nên 
		 \allowdisplaybreaks
		\begin{align*}
			 &a_0+8a_1=2a_2+1
			\\ 
			\Leftrightarrow& \mathrm{C}_n^0+16\mathrm{C}_n^1=8\mathrm{C}_n^2+1 
			 \\
			\Leftrightarrow&  1+16n=\dfrac{8n(n-1)}{2!}+1
			\\ 
			\Leftrightarrow&  n=5.
		\end{align*}
		Suy ra ta  có khai triển $(1+2x)^5=\sum\limits_{k=0}^5\mathrm{C}_5^k 2^k x^k$.\\ Hệ số của khai triển là $a_k=\mathrm{C}_5^k 2^k$.\\
		Ta có 
		\allowdisplaybreaks 
		\begin{align*}
			&a_k ~\text{là  hệ số lớn nhất}
			 \\
			\Leftrightarrow&\heva{&a_k  \geq a_{k+1} \\&a_k \geq a_{k-1}}
			\\ 
			\Leftrightarrow&\heva{&C_5^k 2^k \geq C_5^{k+1} 2^{k+1}  \\
			&C_5^k 2^k \geq  C_5^{k-1} 2^{k-1}
			} 
			 \\
			\Leftrightarrow&\heva{&\dfrac{5!}{k!(5-k)!} 2^k \geq  \dfrac{5!}{(k+1)!(5-k-1)!} 2^{k+1} \\&\dfrac{5!}{k!(5-k)!} 2^k \geq \dfrac{5!}{(k-1)!(5-k+1)!} 2^{k-1}}
			 \\
			\Leftrightarrow&\heva{&\dfrac{1}{5-k} \geq  \dfrac{2}{k+1} \\
			 &\dfrac{2}{k} \geq \dfrac{1}{5-k+1}}
			 \\
			\Leftrightarrow&\heva{&k+1\geq  10-2k \\ &12-2k \geq k}
			\\ 
			\Leftrightarrow& 11\leq 3k  \leq 12\\
			\Leftrightarrow& \dfrac{11}{3} \leq k \leq  4\\
			 \Rightarrow&\hoac{&k=3\\
			&k=4.} 
		\end{align*}
		Vậy hệ số lớn nhất  của khai triển là $a_3=\mathrm{C}_5^32^3=80=a_4=\mathrm{C}_5^42^4=80$.
	} 
\end{ex}


%%%%%%%%%%%%%%%% ĐÚNG/SAI %%%%%%%%%%%%%%%%

% nguồn: MaHoaNH_DungSai_Lop10
\begin{ex}%[0D8H3-4]%[ma:effac21a]
	Trong khai triển $(a+b)^n$,  số hạng tổng quát của khai triển là
	\choiceTF 
		{\True  $\mathrm{C}_n^k{a^{n-k}}{b^k}$}
		 {$\mathrm{C}_n^{k-1}{a^{n+1}}{b^{n-k+1}}$}
		 {\True $\mathrm{C}_n^k{b^{n-k}}{a^k}$}
		 {$\mathrm{C}_n^{k+1}{a^{n-k+1}}{b^{k+1}}$}
	\loigiai{ 
		Theo công thức khai  triển nhị thức Newton ta có $$(a+b)^n=\sum\limits_{k=0}^n\mathrm{C}_n^{k}a^{n-k}b^{k},(0\leq k\leq n).$$
		Hoặc $$(b+a)^n=\sum\limits_{k=0}^n\mathrm{C}_n^{k}a^{k}b^{n-k},(0\leq  k\leq n).$$
		Do đó số hạng tổng quát của khai triển  là $\mathrm{C}_n^{k}a^{n-k}b^{k},(0\leq k\leq n)$ hoặc $\mathrm{C}_n^{k}a^{k}b^{n-k},(0\leq k\leq n)$.
	} 
\end{ex}

% nguồn: MaHoaNH_DungSai_Lop10
\begin{ex}%[0D8H3-4]%[ma:0f1753e2]
	Xét khai triển nhị thức $(a+b)^5$. Xét tính đúng,  sai của mỗi khẳng định sau đây.
	\choiceTF 
		{\True Tổng các số mũ của $a$ và $b$ trong mỗi số hạng bằng  $5$}
		{\True Các hệ số của các cặp số hạng cách đều số hạng đầu và cuối thì  bằng nhau}
		 {Số số hạng của vế phải trong khai triển bằng $5$}
		{\True Số hạng thứ $4$ trong khai triển viết theo bậc  tăng dần của $b$ là $10a^2 b^3$}
	 \loigiai{
		Khai  triển $(a+b)^5=a^5+5a^4b+10a^3{b^2}+10a^2{b^3}+5a{b^4}+b^5$, ta có
		\begin{itemchoice}
			\itemch {\bf Đúng.} Tổng số mũ của $a$ và $b$ trong mỗi số hạng bằng $5$. 
			\itemch {\bf Đúng.} Các hệ số của cặp số hạng cách đều số hạng đầu và cuối thì  bằng nhau.
			\itemch {\bf Sai.} Số số hạng  của vế phải trong khai triển là $6$.
			\itemch {\bf Đúng.} Số hạng thứ $4$ trong khai triển viết  theo bậc tăng dần của $b$ là $\mathrm{C}_5^3 a^{5-3}\cdot b^3=10a^2b^3$.
		\end{itemchoice}
	 }
\end{ex}

% nguồn: MaHoaNH_DungSai_Lop10
\begin{ex}%[0D8H3-4]%[ma:b15c32b4]
	Chọn khẳng định đúng trong khai  triển nhị thức $\left(2x-\dfrac{1}{x^2}\right)^{6}$ với $x\neq 0$.
	\choiceTF 
		 {\True Hệ số của số hạng chứa $x^6$ là $64$}
		 {Hệ số của số hạng chứa $x^3$ là $192$}
		{\True Hệ số của số hạng chứa không chứa  $x$ là $240$}
		{Có số hạng $\dfrac{30}{x^6}$ trong khai triển} 
	 \loigiai{
		Ta có:\begin{eqnarray*}
		\left(2x-\dfrac{1}{x^2}\right)^{6}&=&\sum\limits_{k=0}^{6} \mathrm{C}_{6}^{k} \cdot (2x)^{6-k} \cdot \left(\dfrac{-1}{x^2}\right)^{k}\\ 
		&=&\sum\limits_{k=0}^{6} \mathrm{C}_{6}^{k} \cdot 2^{6-k}\cdot x^{6-k} \cdot  (-1)^k \cdot x^{-2k}\\
		 &=&64 x^{6} - 192 x^{3} + 240 - \dfrac{160}{x^{3}} + \dfrac{60}{x^{6}} - \dfrac{12}{x^{9}} + \dfrac{1}{x^{12}}
	\end{eqnarray*}
	Vì có số hạng không chứa $x$ nên  $6-3k=0 \Leftrightarrow k=2$.\\
	Số  hạng chứa $x^6$ có hệ số là $\mathrm{C}_{6}^{2} \cdot 2^4 \cdot(-1)^{2}= 240$.\\
 }
\end{ex}

% nguồn: MaHoaNH_DungSai_Lop10
\begin{ex}%[0D8H3-4]%[ma:7b0334b2]
	Khai triển  $(x-3)^4$.
	 \choiceTF
		{\True Hệ số của $x^4$ trong khai triển  là $1$}
		{\True Hệ số của $x^3$  trong khai triển là $-12$}
		{\True Hệ  số của $x^2$ trong khai triển là $54$}
		{Tổng các hệ số của  các hạng tử có bậc chẵn trong khai triển bằng $134$}
	\loigiai{ 
	Ta có $(x-3)^4 = x^4 - 12x^3 + 54x^2 - 108x +  81$.
	\begin{itemchoice}
		 \itemch
		 \itemch
		\itemch 
		\itemch Tổng các hệ  số của các hạng tử có bậc chẵn trong khai triển bằng $1+54+81 = 136$.
	\end{itemchoice}
	} 
\end{ex}

% nguồn: MaHoaNH_DungSai_Lop10
\begin{ex}%[0D8H3-4]%[ma:795b650f]
	Xét  khai triển biểu thức $P(x) = (x+1)^4 - (x-1)^4$.
	\choiceTF 
		{\True  Số hạng tự do trong khai triển bằng $0$}
		{Hệ số của $x$ trong  khai triển là $-8$}
		{Tổng các hệ số  trong khai triển bằng $14$}
		 {\True $P(x) > 0$ với mọi $x>0$}
	 \loigiai{
	Ta có  $(x+1)^4 = x^4 + 4x^3 + 6x^2 + 4x + 1$ và $(x-1)^4 = x^4 - 4x^3 + 6x^2 - 4x + 1$ nên $P(x) = 8x^3 + 8x$.
	\begin{itemchoice}
		\itemch 
		\itemch 
		\itemch  Tổng các hệ số trong khai triển là $8+8 = 16$.
		\itemch Ta có $P(x) = 8x(x^2+1) >0  $ với mọi $x>0$.
	\end{itemchoice}
	} 
\end{ex}

% nguồn: MaHoaNH_DungSai_Lop10
\begin{ex}%[0D8H3-4]%[ma:8c429f3b]
	Cho nhị  thức $(x+2)^5$, xét tính đúng sai của các mệnh đề sau.
	 \choiceTF
		{\True Khai triển  $(x+2)^5$ có $6$ số hạng}
		{Hệ  số của số hạng chứa $x^2$ là $40$}
		{\True Tổng hệ số trong khai triển  là $243$}
		{Hệ số của số hạng không chứa  $x$ là $1$}
	 \loigiai{
		\begin{itemchoice}
			\itemch Đúng. Trong khai triển $(a+b)^n$ có $n+1$ số hạng.  Vậy khai triển $\left(x+2\right)^5$ có 6 số hạng.
			\itemch Sai. Vì 
			$$(x+2)^5=\mathrm{C}_5^0x^5+\mathrm{C}_5^1x^4\cdot 2^1+\mathrm{C}_5^2\cdot  x^3\cdot 2^2+\mathrm{C}_5^3 x^2\cdot 2^3+\mathrm{C}_5^4 x^4\cdot 2^4+\mathrm{C}_5^5 2^5=x^5+10x^4+40x^3+80x^2+80x+32$$
			\itemch Đúng. Ta có 
			$$a_0+a_1+a_2+a_3+a_4+a_5=\mathrm{C}_5^0+\mathrm{C}_5^1\cdot 2^1+\mathrm{C}_5^2\cdot 2^2+\mathrm{C}_5^3 \cdot  2^3+\mathrm{C}_5^4\cdot 2^4+\mathrm{C}_5^5=243$$
			\itemch Sai. Hệ số của số hạng không chứa $x$ là  $\mathrm{C}_5^5\cdot 2^5$.
		\end{itemchoice}
	} 
\end{ex}

% nguồn: MaHoaNH_DungSai_Lop10
\begin{ex}%[0D8H3-4]%[ma:3d7e881b]
	Khai  triển $\left(x+\dfrac{1}{x}\right)^{4}$.
	\choiceTF 
		{Hệ số của $x^2$ là  $\dfrac{1}{4}$}
		{\True Số  hạng không chứa $x$ là $6$}
		{\True Hệ số của $x^4$ là  $1$}
		{\True Sau khi khai triển, biểu  thức có $5$ số hạng}
	\loigiai{Ta có 
		\begin{align*}
			\left(x+\dfrac{1}{x}\right)^4&= \mathrm{C}_4^0x^4 + \mathrm{C}_4^1x^3\left(\dfrac{1}{x}\right) + \mathrm{C}_4^2x^2\left(\dfrac{1}{x}\right)^2 + \mathrm{C}_4^3x\left(\dfrac{1}{x}\right)^3 + \mathrm{C}_4^4\left(\dfrac{1}{x}\right)^4\\&= x^4 +  4x^2 + 6 + \dfrac{4}{x^2} + \dfrac{1}{x^4}.
		\end{align*}
		\begin{itemchoice}
			\itemch Hệ số của $x^2$ trong khai triển là  $4$.
			 \itemch Số hạng không chứa $x$ là $6$.
			\itemch Hệ số của $x^4$ trong khai triển là  $1$.
			\itemch Sau  khi khai triển, biểu thức có $5$ số hạng.
		\end{itemchoice}
	 }
\end{ex}

% nguồn: MaHoaNH_DungSai_Lop10
\begin{ex}%[0D8H3-4]%[ma:72e12665]
	Khai triển $P=(x-\sqrt{3})^5$. 
	\choiceTF 
		{Hệ số của  $x^4$ trong khai triển là $5\sqrt{3}$}
		{\True Hệ số  của $x^2$ trong khai triển là $-30\sqrt{3}$}
		{\True  Hệ số của $x^3$ trong khai triển là $30$}
		{\True Hệ số của $x$ trong khai triển  là $45$}
	\loigiai{Ta  có
		\begin{align*}
			(x-\sqrt{3})^5 &=  \mathrm{C}_5^0x^5 + \mathrm{C}_5^1x^4\left(-\sqrt{3}\right) + \mathrm{C}_5^2x^3\left(-\sqrt{3}\right)^2 + \mathrm{C}_5^3x^2\left(-\sqrt{3}\right)^3 + \mathrm{C}_5^4x\left(-\sqrt{3}\right)^4 + \mathrm{C}_5^5\left(-\sqrt{3}\right)^5
			\\&= x^5 - 5\sqrt{3}x^4 + 30x^3 - 30\sqrt{3}x^2 +  45x - 9\sqrt{3}.
		\end{align*}
		\begin{itemchoice}
			\itemch Hệ  số của $x^4$ trong khai triển là $-5\sqrt{3}$.
			 \itemch Hệ số của $x^2$ trong khai triển là $-30\sqrt{3}$.
			\itemch Hệ số của $x^3$ trong khai triển là $30$. 
			\itemch Hệ số của $x$ trong khai triển  là $45$.
		\end{itemchoice}
	} 
\end{ex}

% nguồn: MaHoaNH_DungSai_Lop10
\begin{ex}%[0D8H3-4]%[ma:7ac862f7]
	Khai triển  $Q=(xy-1)^5$.
	 \choiceTF
		{\True Số  hạng có chứa $x^2y^2$ là $-10x^2y^2$}
		{\True Hệ số của $x^4y^4$  trong khai triển là $-5$}
		{\True Hệ số của $x^3y^3$ trong  khai triển là $10$}
		{Hệ số của $xy$ trong khai tri $-5\sqrt{3}$ ển là  $-10$}
	\loigiai{Ta có 
		\begin{align*}
			(xy-1)^5 &= \mathrm{C}_5^0\left(xy\right)^5 + \mathrm{C}_5^1\left(xy\right)^4\left(-1\right) + \mathrm{C}_5^2\left(xy\right)^3\left(-1\right)^2 + \mathrm{C}_5^3\left(xy\right)^2\left(-1\right)^3 + \mathrm{C}_5^4\left(xy\right) \left(-1\right)^4 + \mathrm{C}_5^5\left(-1\right)^5 
			\\&= x^5y^5  - 5x^4y^4 + 10x^3y^3 - 10x^2y^2 + 5xy - 1.
		\end{align*}
		\begin{itemchoice}
			\itemch Số hạng có chứa  $x^2y^2$ trong khai triển là $-10x^2y^2$.
			\itemch Hệ số của $x^4y^4$ trong khai  triển là $-5$.
			\itemch Hệ số của $x^3y^3$ trong  khai triển là $10$.
			\itemch Hệ số của $xy$ trong khai triển là  $5$.
		\end{itemchoice}
	 }
\end{ex}

% nguồn: MaHoaNH_DungSai_Lop10
\begin{ex}%[0D8H3-4]%[ma:441708a8]
	Khai  triển $\left(x+\sqrt{2}\right)^4$. Các mệnh đề sau đúng hay sai?
	 \choiceTF
		{\True Hệ số của  $x^2$ là $12$}
		{Khai triển  trên có $4$ số hạng}
		{\True Hệ số lớn nhất của khai triển là $8\sqrt{2}$} 
		{\True Số hạng không chứa $x$ trong khai triển trên bằng  $4$}
	 \loigiai{
		Ta có  \\
		$\begin{aligned}
		\left(x+\sqrt{2}\right)^4&=\mathrm{C}_4^0 x^4+\mathrm{C}_4^1 x^3\left(\sqrt{2}\right)+\mathrm{C}_4^2  x^2\left(\sqrt{2}\right)^2+\mathrm{C}_4^3 x\left(\sqrt{2}\right)^3+\mathrm{C}_4^4\left(\sqrt{2}\right)^4 \\
		&=x^4+4 \sqrt{2} x^3+12 x^2+8  \sqrt{2} x+4.
	\end{aligned}$
	\begin{itemchoice}
		 \itemch Hệ số của $x^2$ là $12$.
		\itemch Khai triển trên có $5$ số  hạng.
		\itemch Hệ số lớn  nhất của khai triển là $8\sqrt{2}$.
		\itemch Số hạng không chứa $x$ trong khai triển  trên bằng $4$.
	\end{itemchoice}
 }
\end{ex}

% nguồn: MaHoaNH_DungSai_Lop10
\begin{ex}%[0D8H3-4]%[ma:53639077]
	Khai triển  $\left(x+\sqrt{2}\right)^{4}$.
	 \choiceTF
		{\True Hệ số của $x^2$ là  $12$}
		{Hệ số của $x^3$ là  $6\sqrt{2}$}
		{\True Hệ số của $x$  là $8\sqrt{2}$}
		{\True Số hạng không chứa $x$ trong  khai triển trên bằng $4$}
	 \loigiai{Ta có
		\begin{align*}
			 \left(x+\sqrt{2}\right)^4
			&= \mathrm{C}_4^0x^4  + \mathrm{C}_4^1x^3\left(\sqrt{2}\right) + \mathrm{C}_4^2x^2\left(\sqrt{2}\right)^2 + \mathrm{C}_4^3x\left(\sqrt{2}\right)^3 + \mathrm{C}_4^4\left(\sqrt{2}\right)^4
			\\&=  x^4 + 4\sqrt{2}x^3 + 12x^2 + 8\sqrt{2}x + 4.
		\end{align*}
		\begin{itemchoice}
			\itemch Hệ số của  $x^2$ trong khai triển là $12$.
			\itemch Hệ số của $x^3$ trong  khai triển là $4\sqrt{2}$.
			\itemch Hệ số của  $x$ trong khai triển là $8\sqrt{2}$.
			\itemch Số hạng không chứa $x$ trong khai triển trên  bằng $4$.
		\end{itemchoice}
	 }
\end{ex}

% nguồn: MaHoaNH_DungSai_Lop10
\begin{ex}%[0D8H3-4]%[ma:20eb544d]
	Cho khai triển đa thức $P(x)=(2-x)^{12}=a_0+a_1x+a_2x^2+\cdots+a_{11}x^{11}+a_{12} x^{12}$. Xét tính đúng, sai của  mỗi khẳng định sau.
	\choiceTF 
		{\True $a_0  = 2^{12}$}
		{$a_{12}  =-1$}
		 {\True $a_0 + a_1+ \cdots +a_{12} =1$}
		{\True $a_8 =  7920$}
	\loigiai{ 
		Số hạng tổng quát của khai  triển là $\mathrm C_{12}^k 2^{12-k} (-1)^kx^k$, với $k=1,2,\ldots,12$.
		\begin{itemize}
			\item Với $k=0$, ta có $a_0  = 2^{12}$.
			\item Với $k=12$, ta  có $a_{12} = 1$
			\item Thay $x=1$, ta  nhận được $a_0 + a_1+ \cdots +a_{12} =1$.
			\item Với $k=8$,  ta có $a_8=\mathrm C_{12}^82^{4}\cdot(-1)^8=7920$.
		\end{itemize}
		\begin{itemchoice}
			 \itemch Đúng
			\itemch Sai 
			\itemch  Đúng
			 \itemch Đúng
		\end{itemchoice}
	 }
\end{ex}

% nguồn: MaHoaNH_DungSai_Lop10
\begin{ex}%[0D8H3-4]%[ma:f5397aaf]
	 Khai triển $(x+1)^5$.
	 \choiceTF
		{\True Hệ số của $x^4$ là $5$} 
		{\True Số hạng  không chứa $x$ là $1$}
		{$\mathrm{C}_5^0+\mathrm{C}_5^1+\mathrm{C}_5^2+\mathrm{C}_5^3+\mathrm{C}_5^4+\mathrm{C}_5^5=3^5$} 
		{\True  $32\mathrm{C}_5^0+16\mathrm{C}_5^1+8\mathrm{C}_5^2+4\mathrm{C}_5^3+2\mathrm{C}_5^4+\mathrm{C}_5^5=3^5$}
	\loigiai{Ta  có
		\begin{align*}
			(x+1)^5 &= \mathrm{C}_5^0x^5 +  \mathrm{C}_5^1x^4 + \mathrm{C}_5^2x^3 + \mathrm{C}_5^3x^2 + \mathrm{C}_5^4x + \mathrm{C}_5^5\qquad(*)
			\\&= 1 + 5x  + 10x^2 + 10x^3 + 5x^4 + x^5.
		\end{align*}
		\begin{itemchoice}
			\itemch Hệ số của $x^4$  là $5$.
			\itemch Số hạng không chứa  $x$ là $1$
			\itemch Từ khai triển  $(*)$, thay $x=1$, ta được
			\begin{align*}
				(1+1)^5 &= \mathrm{C}_5^01^5 + \mathrm{C}_5^11^4 + \mathrm{C}_5^21^3 + \mathrm{C}_5^31^2 + \mathrm{C}_5^41^1 + \mathrm{C}_5^51^0 \\&= \mathrm{C}_5^0  + \mathrm{C}_5^1 + \mathrm{C}_5^2 + \mathrm{C}_5^3 + \mathrm{C}_5^4 + \mathrm{C}_5^5.
			\end{align*}
			Vậy $\mathrm{C}_5^0 + \mathrm{C}_5^1 + \mathrm{C}_5^2 +  \mathrm{C}_5^3 + \mathrm{C}_5^4 + \mathrm{C}_5^5 = 2^5$.
			\itemch Từ khai triển $(*)$, thay $x=2$,  ta được
			\begin{align*}
				(2+1)^5  &= \mathrm{C}_5^02^5 + \mathrm{C}_5^12^4 + \mathrm{C}_5^22^3 + \mathrm{C}_5^32^2 + \mathrm{C}_5^42^1 + \mathrm{C}_5^52^0\\&= 32\mathrm{C}_5^0 + 16\mathrm{C}_5^1 + 8\mathrm{C}_5^2 + 4\mathrm{C}_5^3 + 2\mathrm{C}_5^4 + \mathrm{C}_5^5 = 3^5.
			\end{align*}
		\end{itemchoice}
	} 
\end{ex}

% nguồn: MaHoaNH_DungSai_Lop10
\begin{ex}%[0D8H3-4]%[ma:4280a760]
	Cho khai triển $\left(1+2x\right)^n=a_0+a_1x+a_2x^2+\ldots+a_nx^n$ (với $n$ là số nguyên dương)  thỏa mãn $a_0+8a_1=2a_2+1$.
	\choiceTF 
		{\True  $a_0=1$}
		{\True $n=5$} 
		{\True Hệ số của $x^5$ trong khai  triển là $32$}
		{Hệ số lớn nhất trong khai triển là  $40$}
	 \loigiai{
		\begin{itemchoice}
			\itemch Ta có $\left(1+2x\right)^n = \displaystyle\sum\limits_{k=0}^{n} \mathrm{C}_n^k \left(2x\right)^k = \displaystyle\sum\limits_{k=0}^{n} \mathrm{C}_n^k 2^k x^k$  ($k \in \mathbb{N}$). \\
			Suy ra $a_k = 2^k \mathrm{C}_n^k$ nên $a_0 = 2^0 \mathrm{C}_n^0 = 1$. 
			\itemch Thay  $a_0 = 1, a_1 = 2\mathrm{C}_n^1, a_2 = 4\mathrm{C}_n^2$ vào giả thiết ta có
			\begin{align*}
				a_0+8a_1=2a_2+1 
				&\Leftrightarrow1 + 16\mathrm{C}_n^1 = 8\mathrm{C}_n^2 + 1\Leftrightarrow 2\mathrm{C}_n^1 =  \mathrm{C}_n^2
				\\&\Leftrightarrow  2\dfrac{n!}{\left(n-1\right) !1!} = \dfrac{n!}{\left(n-2\right) !2!}
				\\&\Leftrightarrow  2n = \dfrac{n\left(n-1\right)}{2}
				\\&\Leftrightarrow n^2 - 5n  = 0 \Leftrightarrow \heva{&n=0
				\\&  n=5.}
			\end{align*}
			Do  $n$ là số nguyên dương nên $n=5$.
			 \itemch
			\begin{align*}
				\left(1+2x\right)^5 &= \mathrm{C}_5^0 + \mathrm{C}_5^1(2x) + \mathrm{C}_5^2 \left(2x\right)^2 + \mathrm{C}_5^3 \left(2x\right)^3 + \mathrm{C}_5^4 \left(2x\right)^4 + \mathrm{C}_5^5 \left(2x\right)^5 \\&=1+10x + 40x^2  + 80x^3 + 80x^4 + 32x^5
			\end{align*}
			\itemch Hệ số lớn nhất trong  khai triển là $80$.
		\end{itemchoice}
	} 
\end{ex}

% nguồn: MaHoaNH_DungSai_Lop10
\begin{ex}%[0D8H3-4]%[ma:ace884e5]
	Cho hai đa thức $f(x)=(2x+1)^{5}$  và $g(x)=(2-3x)^{4}$.
	 \choiceTF
		{Số hạng tự do trong  khai triển của $f(x)$ là $4$}
		{Hệ số  của $x$ trong khai triển của $g(x)$ là $96$}
		{Tổng các hệ  số của các hạng tử trong khai triển $g(x)$ là $-1$}
		{Hệ số  của $x^3$ trong khai triển của $h(x)=f(x)-g(x)$ là $-136$}
	\loigiai{ 
		Ta có 
		\begin{eqnarray*}
			(2x+1)^5 &=& \mathrm{C}_5^0\cdot(2x)^5+\mathrm{C}_5^1\cdot(2x)^4\cdot1^1+\mathrm{C}_5^2\cdot(2x)^3\cdot1^2+\mathrm{C}_5^3\cdot(2x)^2\cdot1^3+\mathrm{C}_5^4\cdot(2x)\cdot1^4+\mathrm{C}_5^5\cdot1^5 \\ 
			&=& 32x^5+80x^4+80x^3+40x^2+10x+1. 
		\end{eqnarray*}
		Lại có 
		\begin{eqnarray*}
			(2-3x)^4 &=&\mathrm{C}_4^0\cdot2^4\cdot(-3x)^0+\mathrm{C}_4^1\cdot2^3\cdot(-3x)^1+\mathrm{C}_4^2\cdot2^2\cdot(-3x)^2+\mathrm{C}_4^3\cdot2^1\cdot(-3x)^3+\mathrm{C}_4^4\cdot2^0\cdot(-3x)^4 \\ 
			&=&  16-96x+216x^2-216x^3+81x^4\\
			&=&  81x^4-216x^3+216x^2-96x+16.
		\end{eqnarray*}
		\begin{itemchoice}
			\itemch Từ khai triển trên, hệ số tự do trong khai triển của $f(x)$  là $1$.
			\itemch Từ khai triển trên, hệ số của $x$ trong khai triển của $g(x)$  là $-96$.
			\itemch Thay $x=1$ vào $g(x)$, ta có tổng các hệ số của các hạng tử trong khai triển  $g(x)$ là
			\[ 
			g(1)=(2-3)^4=(-1)^4=1. 
			\] 
			\itemch Hệ số của $x^3$ trong khai triển của  $h(x)=f(x)-g(x)$ là $80-(-216)=296$.
		\end{itemchoice}
	 }
\end{ex}

% nguồn: MaHoaNH_DungSai_Lop10
\begin{ex}%[0D8V3-4]%[ma:e4c7f3de]
	Xét khai triển  biểu thức $\left(1+x+x^2\right)^4$ (sắp xếp theo thứ tự mũ giảm dần).
	\choiceTF 
		{\True Số hạng tự  do trong khai triển bằng $1$}
		 {\True Trong khai triển có $9$ hạng tử}
		{Số  hạng chính giữa trong khai triển bằng $18x^4$}
		 {\True Tổng các hệ số trong khai triển bằng $81$}
	\loigiai{ 
	Ta có 
	\begin{eqnarray*}
		&\left(1+x+x^2\right)^4 & = \left[(1+x) +  x^2\right]^4\\
		 & &= (1+x)^4 +4(1+x)^3 x^2 + 6(1+x)^2(x^2)^2 + 4(1+x)(x^2)^3 + (x^2)^4\\
		& &= 1+4x+6x^2+4x^3+x^4 + 4x^2(1+3x+3x^2+x^3) +6x^4(1+2x+x^2) +  4x^6(1+x) + x^8\\
		&  &= x^8 + 4x^7 + 10x^6 + 16x^5 + 19x^4 + 16x^3 + 10x^2 + 4x + 1.
	\end{eqnarray*}
	\begin{itemchoice}
		\itemch 
		 \itemch
		\itemch Số hạng chính giữa trong khai triển  là $19x^4$.
		\itemch Tổng các hệ số trong khai triển  bằng $1+4+10+16+19+16+10+4+1=81$.
	\end{itemchoice}
	 }
\end{ex}


%%%%%%%%%%%%%%%% TRẢ LỜI NGẮN %%%%%%%%%%%%%%%%

% nguồn: MaHoaNH_TraLoiNgan_Lop10
\begin{bt}%[0D8H3-4]%[ma:fda70eef]
	Hệ số của $x^2y^2$ trong khai triển nhị thức  $( x+2y)^4$ là
	\shortans{24} 
	\loigiai{ 
		Số hạng thứ $k+1$ trong khai triển  $( x+2y)^4$ là $\mathrm{C}_4^k x^{4-k}(2y)^k$.\\
		Số hạng chứa $x^2y^2$ tương ứng với $k=2$. Hệ số là $\mathrm{C}_4^2\cdot 2^2=24$. 
	 }
\end{bt}

% nguồn: MaHoaNH_TraLoiNgan_Lop10
\begin{bt}%[0D8H3-4]%[ma:fa2f22aa]
	Hệ số của  số hạng chứa $ x^{15}y^9$ trong khai triển nhị thức $( xy-x^2)^{12}$ là
	 \shortans{-220}
	\loigiai{Ta  có $$
		\left(xy-x^2\right)^{12}=\sum_{k=0}^{12}\mathrm{C}_{12}^k(x y)^{12-k}\cdot (-1)^k\cdot x^{2 k}=\sum_{k=0}^{12}\mathrm{C}_{12}^k  (-1)^k\cdot x^{12+k}\cdot y^{12-k}.
		 $$
		Số hạng chứa $x^{15} y^{19}$  ứng với $k=3$.\\
	Hệ  số của số hạng chứa $x^{15} y^9$ là $-\mathrm{C}_{12}^3=-220$.}
\end{bt}

% nguồn: MaHoaNH_TraLoiNgan_Lop10
\begin{bt}%[0D8H3-4]%[ma:74a7ab43]
	Hệ số của $x$ trong khai nhị thức Newton của  biểu thức $(x+3)^4$ bằng bao nhiêu? \shortans{108}
	\loigiai{ Ta  có
		\begin{eqnarray*}
			(x+3)^4&=&\mathrm{C}_4^0\cdot x^4+ \mathrm{C}_4^1\cdot x^3 \cdot 3+\mathrm{C}_4^2\cdot x^2 \cdot 3^2+\mathrm{C}_4^3\cdot  x \cdot 3^3 +\mathrm{C}_4^4\cdot 3^4\\
			&=&x^4+12x^3+54x^2+108x+81. 
		\end{eqnarray*}
		Vậy hệ số  của $x$ bằng $108$.
	} 
\end{bt}

% nguồn: MaHoaNH_TraLoiNgan_Lop10
\begin{bt}%[0D8H3-4]%[ma:9abc7214]
	Hệ số của số hạng chứa $x^2$ trong khai triển biểu thức  $(3x+4)^4$ là bao nhiêu?.\shortans{864}
	 \loigiai{
		Áp  dụng công thức nhị thức Newton, ta được
		\begin{eqnarray*}
			(3x+4)^4&=& \mathrm{C}_4^0\cdot (3x)^4+\mathrm{C}_4^1\cdot (3x)^3\cdot 4+\mathrm{C}_4^2\cdot (3x)^2\cdot 4^2+\mathrm{C}_4^3\cdot (3x)\cdot 4^3+\mathrm{C}_4^4\cdot  4^4\\
			 &=&81x^4+432x^3+864x^2+768x+256.
		\end{eqnarray*}
	 }
\end{bt}

% nguồn: MaHoaNH_TraLoiNgan_Lop10
\begin{bt}%[0D8H3-4]%[ma:105d20e9]
	Tìm hệ số của $x^3$ trong  khai triển $\left(3x-1\right)^5$.
	\shortans{270} 
	\loigiai{ 
	Ta có  khai triển
	\begin{eqnarray*}
		 (3x-1)^5 &=& \mathrm{C}_5^0 \cdot (3x)^5 + \mathrm{C}_5^1 \cdot (3x)^4 \cdot (-1) + \mathrm{C}_5^2 \cdot (3x)^3 \cdot (-1)^2 +\mathrm{C}_5^3 \cdot (3x)^2 \cdot (-1)^3 + \mathrm{C}_5^4 \cdot (3x) \cdot (-1)^4 + \mathrm{C}_5^5 \cdot (-1)^5\\
		&=& 243x^5 - 405x^4 + 270x^3 -  90x^2 + 15x -1.
	\end{eqnarray*}
	Vậy hệ số của số hạng chứa $x^3$ là $270$. 
	} 
\end{bt}

% nguồn: MaHoaNH_TraLoiNgan_Lop10
\begin{bt}%[0D8H3-4]%[ma:54b17c18]
	 Trong khai triển $(3x+2)^5$ thì hệ số của $x^3$ bằng bao nhiêu?
	 \shortans{1080}
	\loigiai{ 
	Ta có 
	$$\begin{aligned}
	 (3x+2)^5
	&= (3x)^5 + 5 (3x)^4\cdot 2 + 10 (3x)^3 2^2 + 10(3x)^22^3 + 5\cdot 3x\cdot 2^4 + 2^5 \\ 
	&= 243x^5 + 810x^4 + 1080x^3 + 720x^3 + 240x +  32.
\end{aligned}$$
Vậy hệ số  của $x^3$ là $1080$.
 }
\end{bt}

% nguồn: MaHoaNH_TraLoiNgan_Lop10
\begin{bt}%[0D8V3-4]%[ma:3e411e02]
	 Tính tổng các hệ số trong khai triển $(1-2x)^5$.
	\shortans{-1} 
	\loigiai{ 
	Đặt $(1-2x)^5 = a_0 +  a_1x+\cdots +a_5x^5$.\\
	Cho $x=1$, ta có tổng các hệ số $a_0 + a_1 + \cdots + a_5  = (1-2)^5 = -1$.
	} 
\end{bt}

% nguồn: MaHoaNH_TraLoiNgan_Lop10
\begin{bt}%[0D8V3-4]%[ma:6b9e5d84]
	Tìm hệ số  của số hạng chứa $x^3y$ trong khai triển $\left(2xy + \dfrac{3}{y}\right)^5$.
	 \shortans{720}
	 \loigiai{
	Số hạng tổng quát trong khai triển có  dạng
	$$ 
	T_{k+1} = \mathrm{C}^{k}_{5}(2xy)^{5-k}\left(\dfrac{3}{y}\right)^k = \mathrm{C}^{k}_{5}2^{5-k}3^k \cdot x^{5-k}y^{5-2k} \text{ với } k =  \overline{0,1,\ldots,5}.
	 $$
	Số hạng chứa $x^3y$ ứng với $\heva{&5-k = 3\\ &5-2k =  1} \Leftrightarrow k=2$.\\
	Vậy  hệ số của số hạng chứa $x^3y$ là $\mathrm{C}^{2}_{5}2^3 3^2 = 720$.
	} 
\end{bt}

% nguồn: MaHoaNH_TraLoiNgan_Lop10
\begin{bt}%[0D8V3-4]%[ma:ea79eb3e]
	Cho biết hệ  số của $x^2$ trong khai triển $(1+2x)^n$ bằng $180$. Tìm $n$.
	 \shortans{10}
	\loigiai{ 
	Số hạng tổng quát của khai triển $(1+2x)^n$ là 
	$$T_{k+1} = \mathrm{C}^{k}_{n}1^{n-k}(2x)^k = \mathrm{C}^{k}_{n} 2^kx^k \ \text{với}  \ k=\overline{0,1,\ldots, n}.$$
	Do hệ số của $x^2$  bằng $180$ nên
	 $$
	\mathrm{C}^{2}_{n}2^2 = 180 \Leftrightarrow 2n(n-1) =  180 \Leftrightarrow \hoac{&n=10&&\text{ (thỏa mãn)}\\ &n=-9 & &\text{ (loại).}}
	 $$
	Vậy $n=10$. 
	} 
\end{bt}

% nguồn: MaHoaNH_TraLoiNgan_Lop10
\begin{bt}%[0D8V3-4]%[ma:7d737630]
	Tìm hệ sô của số hạng chứa $x^3$ trong khai triển  của đa thức $x(2x+1)^4+(x+2)^5$.
	\shortans{64} 
	\loigiai{ 
		\begin{itemize}
			\item  Ta có $x(2x+1)^4=x\sum\limits_{k=0}^4 \mathrm{C}_4^k2^{4-k}x^{4-k}$. \\
			 Số hạng chứa $x^3$ tương ứng với $4-k=2\Rightarrow k=2$.
			\item Ta có  $(x+2)^5=\sum\limits_{k=0}^5 \mathrm{C}_5^k\cdot 2^kx^{5-k}$. \\
			Số hạng chứa $x^3$ tương  ứng với $5-k=3\Rightarrow k=2$.
		\end{itemize}
		Vậy  số hạng cần tìm là $\mathrm{C}_4^2\cdot 2^2x^3+\mathrm{C}_5^2\cdot 2^2x^3=64x^3$.
	 }
\end{bt}

% nguồn: MaHoaNH_TraLoiNgan_Lop10
\begin{bt}%[0D8V3-4]%[ma:c3be74aa]
	Cho $n$ là số tự nhiên thỏa mãn $\mathrm C_n^0+\mathrm C_n^1+\mathrm  C_n^2+\ldots+\mathrm C_n^{n-1}+\mathrm C_n^n=2^{2007}$. Trong khai triển $P(x)=\left(\sqrt[3]{2}x-\sqrt{3}\right)^n$ thành đa thức có tất cả bao nhiêu số hạng với hệ số nguyên dương?
	\shortans{335} 
	\loigiai{ 
		 Ta có
		$$(1+x)^{2007}=\mathrm C_{2007}^0+\mathrm C_{2007}^1x+\mathrm C_{2007}^2x^2+\ldots+\mathrm C_{2007}^{{2007}-1}x^{{2007}-1}+\mathrm C_{2007}^{2007}x^{2007}.$$ 
		Chọn $x=1$ ta được 
		 $$2^{2007}=\mathrm C_{2007}^0+\mathrm C_{2007}^1+\mathrm C_{2007}^2+\ldots+\mathrm C_{2007}^{{2007}-1}+\mathrm C_{2007}^{2007}.$$
		Hệ số  tổng quát trong khai triển $P(x)$ là $\mathrm C_{2007}^k\left(\sqrt[3]{2}\right)^{2007-k}\cdot\left(\sqrt{3}\right)^k$\\
		Để có  hệ số nguyên dương thì $k$ là số chẵn và $2017-k$ chia hết cho $3$.\\
		Ta suy  ra $2007-k$ phải là số lẻ chia hết cho $3$ nên các số đó là $2007,2001,\ldots,3$.\\
		Số số hạng có  hệ số nguyên dương là $\dfrac{2007-3}{6}+1=335$.
	} 
\end{bt}

% nguồn: MaHoaNH_TraLoiNgan_Lop10
\begin{bt}%[0D8V3-4]%[ma:5eab0c7e]
	Tìm hệ số của $x^{4}$  trong khai triển $(2-3x)^{5}+2(3x+1)^{6}$.
	\shortans{3240} 
	 \loigiai{
		Hệ số của $x^{4}$ trong khai triển cũng là tổng hệ số  của $x^{4}$ trong khai triển của hai biểu thức thành phần.
		\begin{itemize}
			\item Số hạng tổng quát trong khai triển $(2-3x)^{5}$ có dạng $\mathrm{C}_{5}^{k} \cdot 2^{5-k} \cdot (-3x)^{k} = \mathrm{C}_{5}^{k} \cdot  2^{5-k} \cdot (-3)^{k} \cdot x^{k}$.\\
			 Số hạng chứa $x^{4}$ tương ứng với $k=4$.\\
			Vậy hệ số của $x^{4}$ trong khai triển $(2-3x)^{5}$ là  $\mathrm{C}_{5}^{4} \cdot 2^{5-4} \cdot (-3)^{4}= 810$.
			\item Số hạng tổng quát trong khai triển $(3x+1)^{6}$  có dạng $\mathrm{C}_{6}^{h} \cdot (3x)^{6-h} \cdot 1^{h} = \mathrm{C}_{6}^{h} \cdot 3^{6-h} \cdot x^{6-h}$.\\
			Số hạng chứa  $x^{4}$ tương ứng với $6-h=4 \Rightarrow h=2$.\\
			Vậy  hệ số của $x^{4}$ trong khai triển $(3x+1)^{6}$ là $\mathrm{C}_{6}^{2} \cdot 3^{6-2}= 1215$.
		\end{itemize}
		Kết luận: hệ số của $x^{4}$ trong khai triển đã cho là $810+2  \cdot 1215=3240$.
	} 
\end{bt}

% nguồn: MaHoaNH_TraLoiNgan_Lop10
\begin{bt}%[0D8V3-4]%[ma:f94b8cb5]
	Tìm hệ số  của $x^{10}$ trong khải triển thành đa thức của $\left(1+x+x^2+x^3\right)^5$.
	 \shortans{101}
	 \loigiai{
		Ta  có $\left(1+x+x^2+x^3\right)^5
		=\left[(1+x)  \cdot\left(1+x^2\right)\right]^5
		=(1+x)^5\cdot \left(1+x^2\right)^5$.\\ 
		Xét khai triển  $(1+x)^5\cdot \left(1+x^2\right)^5
		=\sum\limits_{k=0}^5 \mathrm{C}_5^k x^k\cdot \sum_{l=0}^5  \mathrm{C}_5^l x^{2l}
		=\sum\limits_{k=0}^5\left(\mathrm{C}_5^k\cdot \sum_{l=0}^5  \mathrm{C}_5^l \cdot x^{k+2l}\right)$. \\
		Số hạng chứa $x^{10}$ tương ứng với $k, l$ thỏa mãn $k+2l=10 \Leftrightarrow  k=10-2l$. \\
		Kết hợp với điều  kiện, ta có hệ:
		$\heva{&k=10-2l \\ &0 \leq k \leq 5\, (k \in \mathbb{N}) \\ &0 \leq l \leq 5\, (l \in \mathbb{N})}  \Leftrightarrow (k, l) \in \{(0; 5),(2; 4),(4; 3)\}$.\\
		Vậy hệ số  của $x^{10}$ bằng tổng các $\mathrm{C}_5^k \cdot \mathrm{C}_5^l$ thỏa mãn
		$\mathrm{C}_5^0\cdot \mathrm{C}_5^5+\mathrm{C}_5^2\cdot  \mathrm{C}_5^4+\mathrm{C}_5^4\cdot \mathrm{C}_5^3=101$.
	 }
\end{bt}

% nguồn: MaHoaNH_TraLoiNgan_Lop10
\begin{bt}%[0D8C3-4]%[ma:ff7a46e6]
	Biết $n$ là số nguyên dương thỏa mãn $\mathrm{C}^{3}_{n} + \mathrm{A}^{2}_{n} = 50$ và hệ số của số hạng chứa $x^8$ trong khai triển của $f(x) = \left(\dfrac{n}{2x} + \dfrac{x}{2}\right)^{2n}$ ($x\ne 0$) có dạng $\dfrac{a}{b}$ với $a$, $b$ nguyên dương và $(a,b)=1$. Tính  $a+b$.
	\shortans{809} 
	\loigiai{ 
	Ta có $\mathrm{C}^{3}_{n} + \mathrm{A}^{2}_{n} = 50 \Leftrightarrow \dfrac{n(n-1)(n-2)}{6} + n(n-1) = 50 \Leftrightarrow n^3  +3n^2 -4n - 300 = 0 \Leftrightarrow n=6$.\\
	Từ đó ta được $f(x) = \left(\dfrac{3}{x}+\dfrac{x}{2}\right)^{12}$ có số hạng  tổng quát là
	$$ 
	T_{k+1} =  \mathrm{C}^{k}_{12}\left(\dfrac{3}{x}\right)^{12-k} \left(\dfrac{x}{2}\right)^k
	= \mathrm{C}^{k}_{12}3^{12-k}2^{-k}x^{2k-12} \text{ với }  k = \overline{0,1,\ldots,12}.
	$$ 
	 Số hạng chứa $x^8$ ứng với $2k-12 = 8 \Leftrightarrow k=10$ nên hệ số của $x^8$ trong khai triển của $f(x)$ là
	 $$
	 \mathrm{C}^{10}_{12}3^2 2^{-10} = \dfrac{297}{512}.
	$$ 
	Vậy $a+b= 297+512  = 809$.
	 }
\end{bt}

% nguồn: MaHoaNH_TraLoiNgan_Lop10
\begin{bt}%[0D8C3-4]%[ma:17d7a12c]
	Tìm hệ số của số hạng chứa $x^2$ trong khai triển của biểu thức $P(x) = (3+x-x^2)^n$ với $n$  là số nguyên dương thỏa mãn $\mathrm{C}^{2}_{n} + \dfrac{\mathrm{A}^{3}_{n}}{n} = 12$.
	 \shortans{-54}
	 \loigiai{
	Ta có 
	$$\begin{aligned}
	& \ \mathrm{C}^{2}_{n} + \dfrac{\mathrm{A}^{3}_{n}}{n} = 12 \\ 
	\Leftrightarrow &  \ \dfrac{n(n-1)}{2} + \dfrac{1}{n} \cdot n(n-1)(n-2) = 12 \\
	\Leftrightarrow & \  3n^2 - 7n - 20 = 0 \\
	\Leftrightarrow & \ \hoac{ & n=4 & & \text{(thỏa mãn)} \\ &n=-\dfrac{5}{3} &  & \text{(loại).}}
\end{aligned}$$
 Khi đó
\begin{eqnarray*}
	 &P(x) = \left(3+x-x^2\right)^4 & = \left[(3+x) - x^2\right]^4\\
	& &= (3+x)^4 -4(3+x)^3 x^2 + 6(3+x)^2(x^2)^2 -  4(3+x)(x^2)^3 + (x^2)^4\\
	& &= x^8 - 4 x^7 -  6 x^6 + 32 x^5 + 19 x^4 - 96 x^3 - 54 x^2 + 108 x + 81.
\end{eqnarray*}
Vậy hệ số của số hạng  chứa $x^2$ là $-54$.
 }
\end{bt}


%%%%%%%%%%%%%%%% TỰ LUẬN %%%%%%%%%%%%%%%%

% nguồn: MaHoaNH_TuLuan_Lop10
\begin{bt}%[0D8H3-4]%[ma:53dba33f]
	Tính  tổng các hệ số trong khai triển $(1-2x)^{5}$.
	 \loigiai{
		Đặt  $P(x) = (1-2x)^{5} = a_{0}+a_{1}x+a_{2}x^{2}+\ldots+a_{5}x^{5}$.\\
		Tổng các hệ số của khai triển $P(x)$ là giá trị  của $P(x)$ tại $x=1$.\\
		Cho  $x=1$ ta có tổng các hệ số $$a_{0}+a_{1}+a_{2}+\ldots+a_{5}=(1-2(1))^{5}=(-1)^{5}=-1.$$
	} 
\end{bt}

% nguồn: MaHoaNH_TuLuan_Lop10
\begin{bt}%[0D8H3-4]%[ma:da2056a2]
	Tìm hệ số của số hạng  chứa $x^{4}$ trong khai triển $(2+3x)^{5}$.
	 \loigiai{
		 Xét khai triển $(2+3x)^{5}$ có số hạng tổng quát là $$T_{k+1}=\mathrm{C}_{5}^{k}2^{5-k}(3x)^{k}=\mathrm{C}_{5}^{k}2^{5-k}3^{k}x^{k}.$$
		Số hạng chứa $x^{4}$ trong khai triển ứng với  giá trị $k$ thỏa mãn $k=4$.\\
		Vậy hệ số của số hạng chứa $x^{4}$ trong khai triển là\\ $$\mathrm{C}_{5}^{4}2^{5-4}3^{4} = 5 \cdot  2 \cdot 81 = 810.$$
	 }
\end{bt}

% nguồn: MaHoaNH_TuLuan_Lop10
\begin{bt}%[0D8H3-4]%[ma:28cd0994]
	 Tìm hệ số của số hạng không chứa $x$ trong khai triển $\left( \dfrac{x}{2}+\dfrac{4}{x}\right)^{4}$ với $x\ne0$.
	 \loigiai{
		Xét khai triển $\left( \dfrac{x}{2}+\dfrac{4}{x}\right)^{4}$ (với $x\ne0$) có số  hạng tổng quát là
		$$T_{k+1}=\mathrm{C}_{4}^{k}\left( \dfrac{x}{2}\right) ^{4-k}\left( \dfrac{4}{x}\right) ^{k} = \mathrm{C}_{4}^{k} \dfrac{x^{4-k}}{2^{4-k}} \dfrac{4^k}{x^k} = \mathrm{C}_{4}^{k}  \dfrac{2^{2k}}{2^{4-k}} x^{4-k-k} = \mathrm{C}_{4}^{k} 2^{2k-(4-k)} x^{4-2k} = \mathrm{C}_{4}^{k} 2^{3k-4} x^{4-2k}.$$
		Số hạng không  chứa $x$ trong khai triển tương ứng với giá trị $k$ thỏa mãn $$4-2k=0 \Leftrightarrow 2k=4 \Leftrightarrow k=2.$$
		Vậy hệ  số của số hạng không chứa $x$ trong khai triển là $$\mathrm{C}_{4}^{2} \cdot 2^{3(2)-4} = \mathrm{C}_{4}^{2} \cdot 2^{6-4} = \mathrm{C}_{4}^{2} \cdot 2^{2} = 6 \cdot 4 = 24.$$
	 }
\end{bt}

% nguồn: MaHoaNH_TuLuan_Lop10
\begin{bt}%[0D8H3-4]%[ma:8e7017a6]
	Tìm số hạng chứa $x$  trong khai triển $(3x-2)^{4}$.
	\loigiai{ 
		Xét khai triển $(3x-2)^{4}$ có số  hạng tổng quát là
		$$T_{k+1}=\mathrm{C}_{4}^{k}(3x)^{4-k}(-2)^{k}=\mathrm{C}_{4}^{k}3^{4-k}(-2)^{k}x^{4-k}.$$ 
		Số hạng chứa $x$ trong khai triển ứng với giá trị $k$ thỏa mãn $4-k=1  \Rightarrow k=3$.\\
		Vậy số hạng chứa $x$ trong khai triển là $\mathrm{C}_{4}^{3}3^{4-3}(-2)^{3}x = 4 \cdot  3 \cdot (-8)x = -96x$.
	 }
\end{bt}

% nguồn: MaHoaNH_TuLuan_Lop10
\begin{bt}%[0D8H3-4]%[ma:59093ae9]
	Tìm hệ số của số hạng chứa $x^{3}$ trong khai triển $\left(  x^{3}+\dfrac{1}{x}\right) ^{5}$ (với $x\ne0$).
	 \loigiai{
		Xét khai triển  $\left( x^{3}+\dfrac{1}{x}\right) ^{5}$ (với $x\ne0$) có số hạng tổng quát là
		$$T_{k+1}=\mathrm{C}_{5}^{k}(x^{3})^{5-k}\left( \dfrac{1}{x}\right)  ^{k} = \mathrm{C}_{5}^{k}x^{3(5-k)}x^{-k} = \mathrm{C}_{5}^{k}x^{15-3k-k} = \mathrm{C}_{5}^{k}x^{15-4k}.$$
		Số hạng chứa $x^{3}$  tương ứng với giá trị $k$ thỏa mãn $15-4k=3 \Leftrightarrow 4k=12 \Leftrightarrow k=3$.\\
		Vậy hệ số của số hạng chứa $x^{3}$ là  $\mathrm{C}_{5}^{3}=10$.
	} 
\end{bt}

% nguồn: MaHoaNH_TuLuan_Lop10
\begin{bt}%[0D8H3-4]%[ma:1affd7bb]
	Tìm  số hạng chứa $x^{3}$ trong khai triển $(2x-1)^{4}$.
	\loigiai{ 
		Xét khai  triển $(2x-1)^{4}$ có số hạng tổng quát là
		$$T_{k+1} =  \mathrm{C}_{4}^{k}(2x)^{4-k}(-1)^{k} = (-1)^{k}\mathrm{C}_{4}^{k}2^{4-k}x^{4-k}.$$
		Số hạng chứa $x^{3}$ trong khai  triển ứng với giá trị $k$ thỏa mãn $4-k=3 \Rightarrow k=1$.\\
		Vậy số hạng chứa $x^{3}$ trong khai triển là  $(-1)^{1}\mathrm{C}_{4}^{1}2^{4-1}x^{3} = -4 \cdot 2^{3}x^{3} = -32x^{3}$.
	 }
\end{bt}

% nguồn: MaHoaNH_TuLuan_Lop10
\begin{bt}%[0D8V3-4]%[ma:2fadb727]
	Tìm số hạng không chứa $x$  trong khai triển $\left( \dfrac{3}{x}+2x\right) ^{4}$ với $x\ne0$.
	\loigiai{ 
		Xét khai triển $\left( \dfrac{3}{x}+2x\right) ^{4}$ (với $x\ne0$) có số hạng tổng quát là 
		$$T_{k+1}=\mathrm{C}_{4}^{k}\left( \dfrac{3}{x}\right) ^{4-k}(2x)^{k} = \mathrm{C}_{4}^{k} 3^{4-k} x^{-(4-k)} 2^k x^k = \mathrm{C}_{4}^{k}  3^{4-k} 2^k x^{-4+k+k} = \mathrm{C}_{4}^{k} 2^{k} 3^{4-k} x^{2k-4}.$$
		Số hạng không chứa $x$ trong khai  triển tương ứng với giá trị $k$ thỏa mãn
		$$2k-4=0 \Leftrightarrow 2k=4 \Leftrightarrow  k=2.$$
		Vậy số hạng không chứa $x$ trong khai triển là $\mathrm{C}_{4}^{2} 2^{2} 3^{4-2} = \mathrm{C}_{4}^{2}  2^{2} 3^{2} = 6 \cdot 4 \cdot 9 = 216$.
	} 
\end{bt}

% nguồn: MaHoaNH_TuLuan_Lop10
\begin{bt}%[0D8V3-4]%[ma:161fec42]
	Tìm số  hạng chứa $\dfrac{1}{x^{2}}$ trong khai triển $\left( 2x-\dfrac{1}{x^{2}}\right) ^{4}$ với $x\ne0$.
	\loigiai{ 
		Xét khai  triển $\left( 2x-\dfrac{1}{x^{2}}\right) ^{4}$ (với $x\ne0$) có số hạng tổng quát là\\
		$$T_{k+1}=\mathrm{C}_{4}^{k}(2x)^{4-k}\left( -\dfrac{1}{x^{2}}\right) ^{k} = \mathrm{C}_{4}^{k} 2^{4-k} x^{4-k} (-1)^k x^{-2k} = (-1)^{k}\mathrm{C}_{4}^{k}2^{4-k}x^{4-3k}.$$ 
		Số hạng chứa $\dfrac{1}{x^{2}}$ (tức là $x^{-2}$) trong khai triển tương ứng với  giá trị $k$ thỏa mãn
		$$4-3k=-2 \Leftrightarrow  3k=6 \Leftrightarrow k=2.$$
		Vậy số hạng chứa $\dfrac{1}{x^{2}}$ trong khai triển là  $$(-1)^{2}\mathrm{C}_{4}^{2}2^{4-2}x^{4-3(2)} = 1 \cdot 6 \cdot 2^2 \cdot x^{-2} = 24x^{-2} = \dfrac{24}{x^{2}}.$$
	} 
\end{bt}

% nguồn: MaHoaNH_TuLuan_Lop10
\begin{bt}%[0D8V3-4]%[ma:45e22dcd]
	Tìm số hạng không chứa $x$ trong khai triển $\left(  2x^{2}-\dfrac{1}{x^{2}}\right)^{4}$ với $x\ne0$.
	 \loigiai{
		Xét số hạng tổng  quát
		\begin{eqnarray*}
			T_{k+1}&=&\mathrm{C}_{4}^{k}(2x^{2})^{4-k}\left( -\dfrac{1}{x^{2}}\right)  ^{k}
			=\mathrm{C}_{4}^{k}2^{4-k}(x^2)^{4-k}(-1)^{k}\dfrac{1}{(x^2)^k}\\ &=& \mathrm{C}_{4}^{k}2^{4-k}x^{8-2k}(-1)^{k}x^{-2k} = (-1)^{k}\mathrm{C}_{4}^{k}2^{4-k}x^{8-4k}. 
		\end{eqnarray*}
		Số hạng không chứa $x$ ứng với  $8-4k=0 \Leftrightarrow 4k=8 \Leftrightarrow k=2$.\\
		Vậy số  hạng không chứa x là $$T_{2+1} = (-1)^{2}\mathrm{C}_{4}^{2}2^{4-2} = 1 \cdot 6 \cdot 2^2 = 24.$$
	} 
\end{bt}

% nguồn: MaHoaNH_TuLuan_Lop10
\begin{bt}%[0D8V3-4]%[ma:2716b604]
	Cho khai triển $(x+2y)^3+(2x-y)^3$. Tìm hệ số của $x^2y$ trong  khai triển.
	\loigiai{ 
		Ta có\\ 
		$\begin{aligned}
		(x+2y)^3+(2x-y)^3&=\mathrm{C}_3^0x^3+\mathrm{C}_3^1x^2(2y)+\mathrm{C}_3^2x(2y)^2+\mathrm{C}_3^3(2y)^3+\mathrm{C}_3^0(2x)^3+\mathrm{C}_3^1(2x)^2(-y)  \\
		&\hspace*{7.5cm} +\mathrm{C}_3^2(2x)(-y)^2+\mathrm{C}_3^3(-y)^3  \\
		 &=x^3+6x^2y+12xy^2+8y^3+8x^3-12x^2y+6xy^2-y^3 \\
		&=9x^3-6x^2y+18xy^2+7y^3. 
	\end{aligned}$\\
	Hệ số của $x^2y$ trong khai triển  là $-6x^2y$.
} 
\end{bt}

% nguồn: MaHoaNH_TuLuan_Lop10
\begin{bt}%[0D8V3-4]%[ma:488e7807]
	Cho $n$ là số nguyên dương thỏa  mãn ${\mathrm{C}_n^1+\mathrm{C}_n^2=15}$. Tìm số hạng không chứa $x$ trong khai triển: $\left(x+\dfrac{2}{x^4}\right)^n$.
	\loigiai{ 
		Điều kiện $n \geq 2, n \in \mathbb{N}^*$. Ta có 
		\begin{eqnarray*}
			 \mathrm{C}_n^1+\mathrm{C}_n^2=15
			&\Leftrightarrow& n+\dfrac{n(n-1)}{2}=15  \\
			 &\Leftrightarrow& n^2+n-30=0 \\
			 &\Leftrightarrow& \hoac{&n=5 \\ &n=-6} \Rightarrow n=5.
		\end{eqnarray*}
		Khi đó $\left(x+\dfrac{2}{x^4}\right)^5=\sum\limits_{k=0}^5 \mathrm{C}_5^k \cdot 2^k x^{5-k} \cdot \left(\dfrac{1}{x^4}\right)^k =\sum\limits_{k=0}^5 \mathrm{C}_5^k \cdot  2^kx^{5-5k}$.\\
		 Số hạng không chứa $x$ tương ứng $5-5k=0 \Leftrightarrow k=1$. \\
		Suy  ra số hạng không chứa $x$ là $\mathrm{C}_5^1 \cdot 2^1=10$.
	 }
\end{bt}

% nguồn: MaHoaNH_TuLuan_Lop10
\begin{bt}%[0D8V3-4]%[ma:d7543242]
	Cho $n$ là số nguyên dương  thỏa mãn $\mathrm{C}_{n}^{1}+\mathrm{C}_{n}^{2}=15$.
	Tìm số hạng không chứa $x$ trong khai triển $\left( x+\dfrac{2}{x^{4}}\right) ^{n}$  với $x \ne 0$.
	 \loigiai{
		Điều kiện  $n\ge2, n\in\mathbb{N}^{*}$.\\
		 Ta có
		\begin{eqnarray*}
			\mathrm{C}_{n}^{1}+\mathrm{C}_{n}^{2}=15 &\Leftrightarrow& n+\dfrac{n(n-1)}{2}=15 
			\Leftrightarrow 2n + n(n-1) =  30\\ &\Leftrightarrow& 2n + n^2 - n = 30 \Leftrightarrow n^2+n-30=0 \Leftrightarrow \hoac{n=5 \\n=-6.}
		\end{eqnarray*}
		Vì $n$ là số nguyên dương  nên $n=5$.
		\newline Suy ra khai triển là $\left(  x+\dfrac{2}{x^{4}}\right) ^{5}$.\\ Số hạng tổng quát
		$$T_{k+1} = \mathrm{C}_{5}^{k}x^{5-k}(\dfrac{2}{x^{4}})^{k} = \mathrm{C}_{5}^{k}x^{5-k}2^{k}x^{-4k} = \mathrm{C}_{5}^{k}2^{k}x^{5-5k}.$$ 
		Số hạng không chứa $x$ tương ứng với $5-5k=0 \Leftrightarrow  k=1$.\\
		Suy  ra số hạng không chứa $x$ là $\mathrm{C}_{5}^{1} \cdot 2^{1} = 5 \cdot 2 = 10$.
	} 
\end{bt}

% nguồn: MaHoaNH_TuLuan_Lop10
\begin{bt}%[0D8V3-4]%[ma:0105177f]
	Tìm hệ số  của $x^{10}$ trong khai triển thành đa thức của $(1+x+x^{2}+x^{3})^{5}$.
	\loigiai{ 
		Ta có $$(1+x+x^{2}+x^{3})^{5}=[(1+x)+x^{2}(1+x)]^{5}=[(1+x)(1+x^{2})]^{5}=(1+x)^{5}(1+x^{2})^{5}.$$ 
		Khai  triển $(1+x)^{5} = \displaystyle\sum\limits_{k=0}^{5}\mathrm{C}_{5}^{k}x^{k}$.\\
		 Khai triển $(1+x^{2})^{5} = \displaystyle\sum\limits_{l=0}^{5}\mathrm{C}_{5}^{l}(x^{2})^{l} = \displaystyle\sum\limits_{l=0}^{5}\mathrm{C}_{5}^{l}x^{2l}$.\\
		Do đó $(1+x)^{5}(1+x^{2})^{5}  = \left(\displaystyle\sum\limits_{k=0}^{5}\mathrm{C}_{5}^{k}x^{k}\right) \left(\displaystyle\sum\limits_{l=0}^{5}\mathrm{C}_{5}^{l}x^{2l}\right) = \displaystyle\sum\limits_{k=0}^{5}\displaystyle\sum\limits_{l=0}^{5}\mathrm{C}_{5}^{k}\mathrm{C}_{5}^{l}x^{k+2l}$.\\
		 Số hạng chứa $x^{10}$ tương ứng với $k,l$ thỏa mãn $k+2l=10$.\\
		Với điều kiện $0 \le k \le 5$,  $0 \le l \le 5$ và $k,l \in \mathbb{N}$.\\
		 Từ $k=10-2l$\\
		$0  \le 10-2l \le 5 \Rightarrow -10 \le -2l \le -5 \Rightarrow 5 \le 2l \le 10 \Rightarrow 2.5 \le l \le 5$.\\
		Vì $l$ là số nguyên, $l  \in \{3, 4, 5\}$.
		\begin{itemize}
			\item Nếu  $l=3 \Rightarrow k=10-2(3)=4$ (thỏa $0 \le 4 \le 5$). Cặp $(k,l)=(4,3)$.
			\item Nếu $l=4 \Rightarrow k=10-2(4)=2$ (thỏa $0  \le 2 \le 5$). Cặp $(k,l)=(2,4)$.
			\item Nếu $l=5 \Rightarrow k=10-2(5)=0$ (thỏa $0 \le 0 \le  5$). Cặp $(k,l)=(0,5)$.
		\end{itemize}
		Vậy hệ số của $x^{10}$ là tổng các hệ số $\mathrm{C}_{5}^{k}\mathrm{C}_{5}^{l}$  ứng với các cặp $(k,l)$ tìm được
		$$\mathrm{C}_{5}^{4}\mathrm{C}_{5}^{3} + \mathrm{C}_{5}^{2}\mathrm{C}_{5}^{4} +  \mathrm{C}_{5}^{0}\mathrm{C}_{5}^{5} = (5)(10) + (10)(5) + (1)(1) = 50 + 50 + 1 = 101.$$
	 }
\end{bt}

% nguồn: MaHoaNH_TuLuan_Lop10
\begin{bt}%[0D8V3-4]%[ma:b8f5bb63]
	Cho khai triển $(1+2x)^{n}=a_{0}+a_{1}x+a_{2}x^{2}+\ldots+a_{n}x^{n}$ thỏa mãn $a_{0}+8a_{1}=2a_{2}+1$. Tìm giá  trị của số nguyên dương $n$.
	 \loigiai{
		Ta có  $(1+2x)^{n}=\displaystyle\sum \limits_{k=0}^{n}\mathrm{C}_{n}^{k}(2x)^{k} = \displaystyle\sum \limits_{k=0}^{n}\mathrm{C}_{n}^{k}2^{k}x^{k}$ (với $k\in\mathbb{N}$, $0 \le k \le n$).\\
		Suy  ra
		\begin{eqnarray*}
			a_{k}&=&\mathrm{C}_{n}^{k}2^{k}. 
			 \\
			a_{0}&=&\mathrm{C}_{n}^{0}2^{0} = 1 \cdot  1 = 1.
			\\ 
			a_{1}&=&{C}_{n}^{1}2^{1} =  n \cdot 2 = 2n.
			\\ 
			a_{2}&=&\mathrm{C}_{n}^{2}2^{2} = \dfrac{n(n-1)}{2}  \cdot 4 = 2n(n-1).
		\end{eqnarray*}
		Thay vào giả thiết $a_{0}+8a_{1}=2a_{2}+1$. 
		\begin{eqnarray*}
			1+8(2n) &=&  2(2n(n-1)) + 1 \Leftrightarrow
			1+16n = 4n(n-1) + 1  \\
			\Leftrightarrow16n &=&  4n^2 - 4n
			\Leftrightarrow 4n^2  - 20n = 0
			 \Leftrightarrow 4n(n-5) = 0.
		\end{eqnarray*}
		Do $n$ là số nguyên dương nên  $n \ne 0$, suy ra $n-5=0 \Leftrightarrow n=5$.
	 }
\end{bt}

% nguồn: MaHoaNH_TuLuan_Lop10
\begin{bt}%[0D8V3-4]%[ma:6d75c08e]
	Tìm số hạng có hệ số nguyên trong khai triển thành đa thức của $\left( \dfrac{3}{2}-\dfrac{2}{3}x^{2}\right) ^{n}$ biết $n$ là số  nguyên dương thỏa mãn $\mathrm{C}_{2n+1}^{0}+\mathrm{C}_{2n+1}^{2}+\mathrm{C}_{2n+1}^{4}+...+\mathrm{C}_{2n+1}^{2n}=1024$.
	 \loigiai{
		Ta  có tính chất tổng các tổ hợp chẵn $$\mathrm{C}_{m}^{0}+\mathrm{C}_{m}^{2}+\mathrm{C}_{m}^{4}+... = 2^{m-1}.$$
		Áp dụng  với $m=2n+1$, ta có $$\mathrm{C}_{2n+1}^{0}+\mathrm{C}_{2n+1}^{2}+\mathrm{C}_{2n+1}^{4}+...+\mathrm{C}_{2n+1}^{2n} = 2^{(2n+1)-1} = 2^{2n}.$$
		Theo  đề bài $2^{2n} = 1024$.\\
		Mà $1024 = 2^{10}$. Suy  ra $2n = 10 \implies n=5$.\\
		Ta cần tìm số hạng có  hệ số nguyên trong khai triển $\left( \dfrac{3}{2}-\dfrac{2}{3}x^{2}\right) ^{5}$.\\
		Số hạng tổng quát của  khai triển là
		\begin{eqnarray*}
			T_{k+1} &=& \mathrm{C}_{5}^{k}\left( \dfrac{3}{2}\right) ^{5-k}(-\dfrac{2}{3}x^{2})^{k} =  \mathrm{C}_{5}^{k}\left( \dfrac{3}{2}\right) ^{5-k}(-1)^{k}\left( \dfrac{2}{3}\right) ^{k}(x^2)^k \\
			&=& \mathrm{C}_{5}^{k}(-1)^{k} \dfrac{3^{5-k}}{2^{5-k}}  \dfrac{2^k}{3^k} x^{2k} = \mathrm{C}_{5}^{k}(-1)^{k} 3^{5-k-k} 2^{k-(5-k)} x^{2k}.
			= \mathrm{C}_{5}^{k}(-1)^{k} 3^{5-2k}  2^{2k-5} x^{2k}.
		\end{eqnarray*}
		Hệ số của số hạng này là $A_k = \mathrm{C}_{5}^{k}(-1)^{k} 3^{5-2k} 2^{2k-5}$.  Ta cần $A_k$ là số nguyên ($0 \le k \le 5$).
		\begin{itemize}
			\item $k=0: A_0 = \mathrm{C}_{5}^{0}(-1)^{0} 3^{5} 2^{-5}  = 1 \cdot 1 \cdot 243 \cdot \dfrac{1}{32} = \dfrac{243}{32}$ (Không nguyên).
			\item $k=1:  A_1 = \mathrm{C}_{5}^{1}(-1)^{1} 3^{3} 2^{-3} = 5 \cdot (-1) \cdot 27 \cdot \dfrac{1}{8} = -\dfrac{135}{8}$ (Không nguyên).
			\item $k=2: A_2 = \mathrm{C}_{5}^{2}(-1)^{2} 3^{1} 2^{-1} = 10 \cdot 1 \cdot 3 \cdot \dfrac{1}{2} = 15$ (Nguyên). 
			\item $k=3: A_3 = \mathrm{C}_{5}^{3}(-1)^{3}  3^{-1} 2^{1} = 10 \cdot (-1) \cdot \dfrac{1}{3} \cdot 2 = -\dfrac{20}{3}$ (Không nguyên).
			\item  $k=4: A_4 = \mathrm{C}_{5}^{4}(-1)^{4} 3^{-3} 2^{3} = 5 \cdot 1 \cdot \dfrac{1}{27} \cdot 8 = \dfrac{40}{27}$ (Không nguyên).
			\item $k=5: A_5 = \mathrm{C}_{5}^{5}(-1)^{5} 3^{-5} 2^{5} = 1 \cdot (-1) \cdot  \dfrac{1}{243} \cdot 32 = -\dfrac{32}{243}$ (Không nguyên).
		\end{itemize}
		Vậy số hạng có  hệ số nguyên duy nhất là khi $k=2$, tương ứng với $A_2 x^{2(2)} = 15x^{4}$.
	} 
\end{bt}

% nguồn: MaHoaNH_TuLuan_Lop10
\begin{bt}%[0D8V3-4]%[ma:7233e0e9]
	Tìm số hạng chứa $x^{2}$ trong  khai triển của biểu thức $P(x)=(3+x-x^{2})^{n}$ với n là số nguyên dương thỏa mãn $\mathrm{C}_{n}^{2}+\dfrac{\mathrm{A}_{n}^{3}}{n}=12$.
	\loigiai{ 
		Điều kiện:  $n\in\mathbb{Z}, n\ge3$ (do có $A_n^3$).\\
		Xét phương  trình
		\begin{eqnarray*}
			&&\mathrm{C}_{n}^{2}+\dfrac{A_{n}^{3}}{n}=12\\ &\Leftrightarrow& 
			\dfrac{n(n-1)}{2!} +  \dfrac{n(n-1)(n-2)}{n} = 12 \Leftrightarrow
			\dfrac{n(n-1)}{2} + (n-1)(n-2) = 12  \\ &\Leftrightarrow&
			n(n-1) +  2(n-1)(n-2) = 24 \Leftrightarrow
			(n-1)[n + 2(n-2)] = 24 \\&\Leftrightarrow& 
			(n-1)(n + 2n -  4) = 24\Leftrightarrow
			(n-1)(3n - 4)  = 24 \\&\Leftrightarrow&
			3n^2  - 4n - 3n + 4 = 24 \Leftrightarrow
			3n^2 - 7n - 20  = 0\\
			&\Leftrightarrow& \hoac{n=4\\n=-\dfrac{5}{3}.} 
		\end{eqnarray*}
		Vậy  $n=4$.\\
		Ta cần tìm số hạng chứa $x^2$ trong khai triển $P(x)=(3+x-x^{2})^{4}  = (3+(x-x^2))^4$.
		$$P(x) = \displaystyle\sum_{k=0}^{4} \mathrm{C}_{4}^{k} 3^{4-k} (x-x^2)^k = \displaystyle\sum_{k=0}^{4} \mathrm{C}_{4}^{k}  3^{4-k} [x(1-x)]^k = \displaystyle\sum_{k=0}^{4} \mathrm{C}_{4}^{k} 3^{4-k} x^k (1-x)^k.$$
		Ta có $(1-x)^k =  \displaystyle\sum\limits_{i=0}^{k} \mathrm{C}_{k}^{i} 1^{k-i} (-x)^i = \displaystyle\sum \limits_{i=0}^{k} \mathrm{C}_{k}^{i} (-1)^i x^i$.\\
		Do đó $P(x) = \displaystyle\sum\limits_{k=0}^{4} \mathrm{C}_{4}^{k} 3^{4-k} x^k \left( \displaystyle\sum\limits_{i=0}^{k}  \mathrm{C}_{k}^{i} (-1)^i x^i \right) = \displaystyle\sum\limits_{k=0}^{4} \displaystyle\sum\limits_{i=0}^{k} \mathrm{C}_{4}^{k} \mathrm{C}_{k}^{i} 3^{4-k} (-1)^i x^{k+i}$.\\
		 Ta cần tìm số hạng chứa $x^2$, nên $k+i=2$.
		Điều kiện $0 \le i \le k \le 4$, $i$, $k \in \mathbb{Z}$.\\ 
		Các cặp $(i,k)$ thỏa mãn\begin{itemize}
		\item Nếu $k=1$, $1+i=2 \Rightarrow i=1$ (thỏa $0 \le 1 \le  1$). Cặp $(1,1)$.
		\item Nếu $k=2$, $2+i=2 \Rightarrow i=0$ (thỏa $0 \le 0  \le 2$). Cặp $(0,2)$.
	\end{itemize}
	Hệ số của $x^2$ là tổng các hệ số tương  ứng với các cặp $(i,k)$ này\\
	Với $(i,k)=(1,1)$, hệ số là $\mathrm{C}_{4}^{1} \mathrm{C}_{1}^{1}  3^{4-1} (-1)^1 = 4 \cdot 1 \cdot 3^3 \cdot (-1) = 4 \cdot 27 \cdot (-1) = -108$.\\
	Với $(i,k)=(0,2)$, hệ số là $\mathrm{C}_{4}^{2} \mathrm{C}_{2}^{0} 3^{4-2} (-1)^0 = 6 \cdot 1 \cdot 3^2 \cdot  1 = 6 \cdot 9 \cdot 1 = 54$.\\
	Tổng hệ số của $x^2$ là $-108 + 54 =  -54$.\\
	Vậy số hạng  chứa $x^2$ là $-54x^2$.
} 
\end{bt}
